%% ============================================================
%% PARTE 7 - ATIVIDADES DE CADA SOM (R203)
%% Atividades por fonema: discriminacao, classificacao e contraste.
%% Nao sao atividades por letra nem por silaba (ja existentes).
%% Gerado por /tmp/opencode/r203_sons_atividades.py (reproduzivel).
%% ============================================================
\chapter{Atividades de Cada Som}

Nas lições e nas folhas Kumon você trabalhou com as \textbf{letras} e com as
\textbf{sílabas}. Estas atividades trabalham com o \textbf{SOM} de cada fone:
você vai \emph{escutar}, \emph{separar} e \emph{comparar} sons, mesmo quando o
mesmo som aparece escrito de jeitos diferentes (por exemplo, o som /k/ em
\textbf{c}asa, \textbf{qu}eijo e \textbf{k}art).

\section{Como usar estas atividades (professor e família)}

\begin{itemize}[leftmargin=1.6cm]
  \item \textbf{Fale devagar} e repita quantas vezes precisar. Ouça antes de escrever.
  \item \textbf{Não é prova}: não há nota certa nem errada. O objetivo é perceber o som.
  \item \textbf{As variações regionais são aceitas}: em algumas regiões TIA se lê
        \emph{tchia} e em outras \emph{tia}; o R forte pode soar como \emph{h}, \emph{x}
        ou \emph{rr}. Todas são formas corretas de falar.
  \item Os recursos de \textbf{mão na garganta} (sentir a vibração) e \textbf{mão no nariz}
        (sentir o ar) são \emph{apoios opcionais}. Nunca force imitação ou contato.
  \item Estas atividades são \textbf{observacionais e lúdicas}. Elas não diagnosticam
        dificuldades de fala ou audição. Se houver preocupação, converse com a família
        e procure um fonoaudiólogo ou psicopedagogo.
  \item Habilidades BNCC trabalhadas: \textbf{EF01LP03, EF01LP05, EF01LP07, EF01LP08,
        EF01LP09}.
\end{itemize}

\section{Quadro-síntese dos sons}

\begin{center}
\begin{tabular}{@{}cp{2.6cm}p{5.2cm}p{5.2cm}@{}}
\hline
\textbf{Nº} & \textbf{Som} & \textbf{Quem faz esse som (grafias)} & \textbf{Exemplos} \\
\hline
1 & \textbf{/a/} & A tônico & \textit{pá, casa, sapato, pata, bala} \\
2 & \textbf{/a/ fraco} & A átono (fraco) & \textit{cama, pata, bala, banana, jaca} \\
3 & \textbf{/é/ aberto} & É aberto & \textit{pé, fé, café, chapéu, bola} \\
4 & \textbf{/e/} & Ê fechado & \textit{mês, você, cedo, mesa, dedo} \\
5 & \textbf{/i/} & I & \textit{ilha, pipa, bico, sino, fivela} \\
6 & \textbf{/ó/ aberto} & Ó aberto & \textit{pó, nó, só, avó, bola, horta} \\
7 & \textbf{/o/} & Ô fechado & \textit{avô, bolo, fogo, coro, gosto} \\
8 & \textbf{/u/} & U & \textit{uva, lua, unha, urubu, chuchu} \\
9 & \textbf{/ã/} & Ã (til) & \textit{anã, maçã, romã, ímã} \\
10 & \textbf{/ã/ fraco} & AN / AM (fraco) & \textit{canta, tampa, banda, campo, andar} \\
11 & \textbf{/e/ nasal} & EN / EM & \textit{tempo, sempre, cento, pente, cem} \\
12 & \textbf{/i/ nasal} & IN / IM & \textit{sim, fim, mim, limpo, pinto} \\
13 & \textbf{/o/ nasal} & ON / OM & \textit{bom, som, tom, ponte, pomba} \\
14 & \textbf{/u/ nasal} & UN / UM & \textit{um, algum, jejum, comum, atum} \\
15 & \textbf{/j/} & I e E semivogais & \textit{pai, vai, cai, boi, papai, herói} \\
16 & \textbf{/w/} & U e L semivogais & \textit{pau, mau, céu, réu; sol, mel, farol; qual, quando} \\
17 & \textbf{/p/} & P & \textit{pato, papai, pipa, pote, pulo, sapo} \\
18 & \textbf{/b/} & B & \textit{bola, baba, bebê, bota, botão} \\
19 & \textbf{/t/} & T & \textit{tatu, teto, tudo, tela, pato, pote} \\
20 & \textbf{/d/} & D & \textit{dado, dedo, doca, duna, dia, cidade} \\
21 & \textbf{/k/} & C, Q, K & \textit{casa, copo, cubo, queijo, quero, aqui, kiwi, kart} \\
22 & \textbf{/g/} & G & \textit{gato, gola, gude; gema, girafa, gelo, página} \\
23 & \textbf{/f/} & F & \textit{faca, foca, figo, fumo, fada, café, sofá} \\
24 & \textbf{/v/} & V & \textit{vaca, vela, vida, vovô, vulcão, ovo, ave} \\
25 & \textbf{/s/} & S, C, Ç, X, Z & \textit{sapo, sopa, sala, sino; cedo, céu, doce; caça, poço; paz, luz; próximo} \\
26 & \textbf{/z/} & Z, S, X & \textit{zero, zebra, zelo; casa, mesa, asa; exame, exemplo} \\
27 & \textbf{/ch/ (x)} & X, CH & \textit{xícara, xale, caixa, peixe; chuva, chave, chão, bicho} \\
28 & \textbf{/j/} & J, G (e, i) & \textit{jaca, jogo, jipe, juba, jejum; gelo, girafa, gema, página} \\
29 & \textbf{/m/} & M & \textit{mala, mão, mim, mar, mel; (fim) bom, som, campo} \\
30 & \textbf{/n/} & N & \textit{nada, nena, nó, nuvem, cana, sino} \\
31 & \textbf{/nh/} & NH & \textit{ninho, unha, sonho, banho, linha, vinho, aranha} \\
32 & \textbf{/lh/} & LH & \textit{milho, alho, filho, olho, ilha, coelho} \\
33 & \textbf{/l/} & L & \textit{lata, lua, leão, sala, bala, lápis} \\
34 & \textbf{/r/ fraco} & R fraco & \textit{cara, pera, aro, barato, coração} \\
35 & \textbf{/h/} & R forte & \textit{rato, rua, rede, rosa, carro; mar, amor, flor} \\
36 & \textbf{/tch/} & T + I (regional) & \textit{tia, leite, noite, tchau, quente, gente (variação regional)} \\
37 & \textbf{/dj/} & D + I (regional) & \textit{dia, dica, dente, cidade, amigo (variação regional)} \\
38 & \textbf{H mudo} & H (sem som) & \textit{hoje (lê-se OJE), hora (ORA), hotel (OTEL), herói (ERÓI)} \\
\hline
\end{tabular}
\end{center}

\section{Atividades de Cada Som}

Cada página a seguir é uma \textbf{plancheta reproduzível} (pode ser copiada para
uso escolar não comercial). Guarde as planchetas feitas: elas mostram o progresso
da criança no reconhecimento de cada som.


\plancheta{Atividade de Som 1/38 --- /a/ --- PÁ, CASA e SAPATO}
\begin{exercicio}[Som /a/ [[a]] --- Atividade de Som (não de letra nem de sílaba)]
\noindent\textbf{Nome:} \underline{\hspace{3.2cm}} \quad \textbf{Data:} \underline{\hspace{1.6cm}}

\noindent\textbf{\faMusic~O som /a/ {[a]}:} pá, casa, sapato, pata, bala.
\newline\textbf{\faPencilRuler~Quem faz esse som:} A tônico. \hfill \textbf{\faCheckCircle~Voz:} com voz.

\vspace{0.4em}
\textbf{A) CAÇA AO SOM} --- Leia em voz alta e compare os sons. Marque a coluna certa para cada palavra.

\begin{center}
\begin{tabular}{@{}ll@{}}
\fbox{\textbf{SIM}} \textit{casa} \quad \fbox{\textbf{NÃO}} \textit{gato} \\
\fbox{\textbf{SIM}} \textit{sapato} \quad \fbox{\textbf{NÃO}} \textit{bola} \\
\fbox{\textbf{SIM}} \textit{pata} \quad \fbox{\textbf{NÃO}} \textit{dedo} \\
\fbox{\textbf{SIM}} \textit{bala} \quad \fbox{\textbf{NÃO}} \textit{bife} \\
\end{tabular}
\end{center}

\textbf{B) SEPARE} --- Escreva as palavras na coluna certa.

\begin{center}
\begin{tabular}{@{}p{4.6cm}|p{4.6cm}@{}}
\hline
\textbf{TEM o som} & \textbf{NÃO tem o som} \\
\hline
\textbf{cama} & \textbf{mesa} \\
\textbf{faca} & \textbf{pente} \\
\textbf{sala} & \textbf{copo} \\
\hline
\end{tabular}
\end{center}

\textbf{C) COMPLETE} --- Complete com a grafia do som e confira entre parênteses.

\begin{center}
\large p\underline{\hspace{0.7cm}}ta (pata) \quad c\underline{\hspace{0.7cm}}sa (casa) \quad b\underline{\hspace{0.7cm}}la (bala) \quad f\underline{\hspace{0.7cm}}ca (faca)
\end{center}

\textbf{D) OUÇA E COMPARE} --- A forte e A fraco.

\begin{center}
\begin{tabular}{@{}p{4.2cm}p{6.2cm}p{1.6cm}@{}}
\hline
\textbf{Palavra} & \textbf{Compare} & \textbf{Marque} \\
\hline
\textbf{casa} & o primeiro A é forte [a], o último é fraco & \textbf{SIM} \\
\textbf{cama} & CA é forte; o A final é fraco & \textbf{SIM} \\
\textbf{pata} & PA é forte; o A final é fraco & \textbf{SIM} \\
\hline
\end{tabular}
\end{center}

\caixapausa{Se precisar, pare aqui. Respire fundo e continue quando quiser.}
\caixapista{Boca bem aberta, maxilar relaxado. Mão na garganta para sentir a vibração da voz.}
\end{exercicio}

\plancheta{Atividade de Som 2/38 --- /a/ fraco --- CAMA, PATA e BALA}
\begin{exercicio}[Som /a/ fraco [[6]] --- Atividade de Som (não de letra nem de sílaba)]
\noindent\textbf{Nome:} \underline{\hspace{3.2cm}} \quad \textbf{Data:} \underline{\hspace{1.6cm}}

\noindent\textbf{\faMusic~O som /a/ fraco {[6]}:} cama, pata, bala, banana, jaca.
\newline\textbf{\faPencilRuler~Quem faz esse som:} A átono (fraco). \hfill \textbf{\faCheckCircle~Voz:} com voz (fraco).

\vspace{0.4em}
\textbf{A) CAÇA AO SOM} --- Leia em voz alta e compare os sons. Marque a coluna certa para cada palavra.

\begin{center}
\begin{tabular}{@{}ll@{}}
\fbox{\textbf{SIM}} \textit{cama} \quad \fbox{\textbf{NÃO}} \textit{pato} \\
\fbox{\textbf{SIM}} \textit{pata} \quad \fbox{\textbf{NÃO}} \textit{bolo} \\
\fbox{\textbf{SIM}} \textit{banana} \quad \fbox{\textbf{NÃO}} \textit{dedo} \\
\fbox{\textbf{SIM}} \textit{bala} \quad \fbox{\textbf{NÃO}} \textit{sapo} \\
\end{tabular}
\end{center}

\textbf{B) SEPARE} --- Escreva as palavras na coluna certa.

\begin{center}
\begin{tabular}{@{}p{4.6cm}|p{4.6cm}@{}}
\hline
\textbf{TEM o som} & \textbf{NÃO tem o som} \\
\hline
\textbf{jaca} & \textbf{sino} \\
\textbf{faca} & \textbf{carro} \\
\textbf{sala} & \textbf{pote} \\
\hline
\end{tabular}
\end{center}

\textbf{C) COMPLETE} --- Complete com a grafia do som e confira entre parênteses.

\begin{center}
\large cam\underline{\hspace{0.7cm}} (cama) \quad pat\underline{\hspace{0.7cm}} (pata) \quad bal\underline{\hspace{0.7cm}} (bala) \quad sal\underline{\hspace{0.7cm}} (sala)
\end{center}

\textbf{D) OUÇA E COMPARE} --- A forte e A fraco.

\begin{center}
\begin{tabular}{@{}p{4.2cm}p{6.2cm}p{1.6cm}@{}}
\hline
\textbf{Palavra} & \textbf{Compare} & \textbf{Marque} \\
\hline
\textbf{casa} & A forte no começo; A fraco no fim & \textbf{SIM} \\
\textbf{cama} & mesma palavra: ouça os dois As & \textbf{SIM} \\
\textbf{banana} & BA forte; NA e NA fracos & \textbf{SIM} \\
\hline
\end{tabular}
\end{center}

\caixapausa{Se precisar, pare aqui. Respire fundo e continue quando quiser.}
\caixapista{Boca quase fechada, som fraquinho. Mão na garganta: vibra, mas fraquinho.}
\end{exercicio}

\plancheta{Atividade de Som 3/38 --- /é/ aberto --- PÉ, FÉ e CHAPÉU}
\begin{exercicio}[Som /é/ aberto [[E]] --- Atividade de Som (não de letra nem de sílaba)]
\noindent\textbf{Nome:} \underline{\hspace{3.2cm}} \quad \textbf{Data:} \underline{\hspace{1.6cm}}

\noindent\textbf{\faMusic~O som /é/ aberto {[E]}:} pé, fé, café, chapéu, bola.
\newline\textbf{\faPencilRuler~Quem faz esse som:} É aberto. \hfill \textbf{\faCheckCircle~Voz:} com voz.

\vspace{0.4em}
\textbf{A) CAÇA AO SOM} --- Leia em voz alta e compare os sons. Marque a coluna certa para cada palavra.

\begin{center}
\begin{tabular}{@{}ll@{}}
\fbox{\textbf{SIM}} \textit{pé} \quad \fbox{\textbf{NÃO}} \textit{dedo} \\
\fbox{\textbf{SIM}} \textit{fé} \quad \fbox{\textbf{NÃO}} \textit{bolo} \\
\fbox{\textbf{SIM}} \textit{café} \quad \fbox{\textbf{NÃO}} \textit{sino} \\
\fbox{\textbf{SIM}} \textit{chapéu} \quad \fbox{\textbf{NÃO}} \textit{lápis} \\
\end{tabular}
\end{center}

\textbf{B) SEPARE} --- Escreva as palavras na coluna certa.

\begin{center}
\begin{tabular}{@{}p{4.6cm}|p{4.6cm}@{}}
\hline
\textbf{TEM o som} & \textbf{NÃO tem o som} \\
\hline
\textbf{avó} & \textbf{chapéu} \\
\textbf{pente} & \textbf{pó} \\
\textbf{fé} & \textbf{você} \\
\hline
\end{tabular}
\end{center}

\textbf{C) COMPLETE} --- Complete com a grafia do som e confira entre parênteses.

\begin{center}
\large p\underline{\hspace{0.7cm}} (pé) \quad f\underline{\hspace{0.7cm}} (fé) \quad caf\underline{\hspace{0.7cm}} (café) \quad ch\underline{\hspace{0.7cm}}péu (chapéu)
\end{center}

\textbf{D) OUÇA E COMPARE} --- É aberto e Ê fechado.

\begin{center}
\begin{tabular}{@{}p{4.2cm}p{6.2cm}p{1.6cm}@{}}
\hline
\textbf{Palavra} & \textbf{Compare} & \textbf{Marque} \\
\hline
\textbf{pé} & abre a boca [e aberto] & \textbf{SIM} \\
\textbf{dedo} & boca mais fechada [e] & \textbf{NÃO} \\
\textbf{café} & abre no fim [e aberto] & \textbf{SIM} \\
\hline
\end{tabular}
\end{center}

\caixapausa{Se precisar, pare aqui. Respire fundo e continue quando quiser.}
\caixapista{Boca mais aberta que no Ê. Mão na frente da boca: sente o ar.}
\end{exercicio}

\plancheta{Atividade de Som 4/38 --- /e/ --- MÊS, VOCÊ e CEDO}
\begin{exercicio}[Som /e/ [[e]] --- Atividade de Som (não de letra nem de sílaba)]
\noindent\textbf{Nome:} \underline{\hspace{3.2cm}} \quad \textbf{Data:} \underline{\hspace{1.6cm}}

\noindent\textbf{\faMusic~O som /e/ {[e]}:} mês, você, cedo, mesa, dedo.
\newline\textbf{\faPencilRuler~Quem faz esse som:} Ê fechado. \hfill \textbf{\faCheckCircle~Voz:} com voz.

\vspace{0.4em}
\textbf{A) CAÇA AO SOM} --- Leia em voz alta e compare os sons. Marque a coluna certa para cada palavra.

\begin{center}
\begin{tabular}{@{}ll@{}}
\fbox{\textbf{SIM}} \textit{mês} \quad \fbox{\textbf{NÃO}} \textit{ela} \\
\fbox{\textbf{SIM}} \textit{você} \quad \fbox{\textbf{NÃO}} \textit{pato} \\
\fbox{\textbf{SIM}} \textit{cedo} \quad \fbox{\textbf{NÃO}} \textit{sapo} \\
\fbox{\textbf{SIM}} \textit{mesa} \quad \fbox{\textbf{NÃO}} \textit{bola} \\
\end{tabular}
\end{center}

\textbf{B) SEPARE} --- Escreva as palavras na coluna certa.

\begin{center}
\begin{tabular}{@{}p{4.6cm}|p{4.6cm}@{}}
\hline
\textbf{TEM o som} & \textbf{NÃO tem o som} \\
\hline
\textbf{dedo} & \textbf{pá} \\
\textbf{mesa} & \textbf{café} \\
\textbf{fivela} & \textbf{cedo} \\
\hline
\end{tabular}
\end{center}

\textbf{C) COMPLETE} --- Complete com a grafia do som e confira entre parênteses.

\begin{center}
\large m\underline{\hspace{0.7cm}}s (mês) \quad c\underline{\hspace{0.7cm}}do (cedo) \quad m\underline{\hspace{0.7cm}}sa (mesa) \quad d\underline{\hspace{0.7cm}}do (dedo)
\end{center}

\textbf{D) OUÇA E COMPARE} --- Ê fechado e É aberto.

\begin{center}
\begin{tabular}{@{}p{4.2cm}p{6.2cm}p{1.6cm}@{}}
\hline
\textbf{Palavra} & \textbf{Compare} & \textbf{Marque} \\
\hline
\textbf{dedo} & fecha [e] & \textbf{NÃO} \\
\textbf{pé} & abre [e aberto] & \textbf{SIM} \\
\textbf{você} & fecha [e] & \textbf{NÃO} \\
\hline
\end{tabular}
\end{center}

\caixapausa{Se precisar, pare aqui. Respire fundo e continue quando quiser.}
\caixapista{Boca entreaberta, lábios levemente esticados. Mão na frente da boca: ar suave.}
\end{exercicio}

\plancheta{Atividade de Som 5/38 --- /i/ --- ILHA, PIPA e BICO}
\begin{exercicio}[Som /i/ [[i]] --- Atividade de Som (não de letra nem de sílaba)]
\noindent\textbf{Nome:} \underline{\hspace{3.2cm}} \quad \textbf{Data:} \underline{\hspace{1.6cm}}

\noindent\textbf{\faMusic~O som /i/ {[i]}:} ilha, pipa, bico, sino, fivela.
\newline\textbf{\faPencilRuler~Quem faz esse som:} I. \hfill \textbf{\faCheckCircle~Voz:} com voz.

\vspace{0.4em}
\textbf{A) CAÇA AO SOM} --- Leia cada par com o dedo no nariz: marca NARIZ para quem faz o nariz vibrar e BOCA para quem não faz.

\begin{center}
\begin{tabular}{@{}ll@{}}
\fbox{\textbf{SIM}} \textit{ilha} \quad \fbox{\textbf{NÃO}} \textit{bola} \\
\fbox{\textbf{SIM}} \textit{pipa} \quad \fbox{\textbf{NÃO}} \textit{dedo} \\
\fbox{\textbf{SIM}} \textit{bico} \quad \fbox{\textbf{NÃO}} \textit{sapo} \\
\fbox{\textbf{SIM}} \textit{sino} \quad \fbox{\textbf{NÃO}} \textit{pato} \\
\end{tabular}
\end{center}

\textbf{B) SEPARE} --- Escreva as palavras na coluna certa.

\begin{center}
\begin{tabular}{@{}p{4.6cm}|p{4.6cm}@{}}
\hline
\textbf{TEM o som} & \textbf{NÃO tem o som} \\
\hline
\textbf{fivela} & \textbf{copo} \\
\textbf{livro} & \textbf{mesa} \\
\textbf{amigo} & \textbf{rato} \\
\hline
\end{tabular}
\end{center}

\textbf{C) COMPLETE} --- Complete com a grafia do som e confira entre parênteses.

\begin{center}
\large \underline{\hspace{0.7cm}}lha (ilha) \quad p\underline{\hspace{0.7cm}}pa (pipa) \quad b\underline{\hspace{0.7cm}}co (bico) \quad s\underline{\hspace{0.7cm}}no (sino)
\end{center}

\textbf{D) OUÇA E COMPARE} --- I oral e I nasal.

\begin{center}
\begin{tabular}{@{}p{4.2cm}p{6.2cm}p{1.6cm}@{}}
\hline
\textbf{Palavra} & \textbf{Compare} & \textbf{Marque} \\
\hline
\textbf{si} & nariz não vibra & \textbf{NÃO} \\
\textbf{sim} & nariz vibra [i nasal] & \textbf{SIM} \\
\textbf{vi} & nariz não vibra & \textbf{NÃO} \\
\hline
\end{tabular}
\end{center}

\caixapausa{Se precisar, pare aqui. Respire fundo e continue quando quiser.}
\caixapista{Lábios esticados em sorriso. Mão no canto da boca para sentir o sorriso.}
\end{exercicio}

\plancheta{Atividade de Som 6/38 --- /ó/ aberto --- PÓ, NÓ e AVÓ}
\begin{exercicio}[Som /ó/ aberto [[O]] --- Atividade de Som (não de letra nem de sílaba)]
\noindent\textbf{Nome:} \underline{\hspace{3.2cm}} \quad \textbf{Data:} \underline{\hspace{1.6cm}}

\noindent\textbf{\faMusic~O som /ó/ aberto {[O]}:} pó, nó, só, avó, bola, horta.
\newline\textbf{\faPencilRuler~Quem faz esse som:} Ó aberto. \hfill \textbf{\faCheckCircle~Voz:} com voz.

\vspace{0.4em}
\textbf{A) CAÇA AO SOM} --- Leia em voz alta e compare os sons. Marque a coluna certa para cada palavra.

\begin{center}
\begin{tabular}{@{}ll@{}}
\fbox{\textbf{SIM}} \textit{pó} \quad \fbox{\textbf{NÃO}} \textit{pé} \\
\fbox{\textbf{SIM}} \textit{nó} \quad \fbox{\textbf{NÃO}} \textit{fé} \\
\fbox{\textbf{SIM}} \textit{só} \quad \fbox{\textbf{NÃO}} \textit{café} \\
\fbox{\textbf{SIM}} \textit{avó} \quad \fbox{\textbf{NÃO}} \textit{chapéu} \\
\end{tabular}
\end{center}

\textbf{B) SEPARE} --- Escreva as palavras na coluna certa.

\begin{center}
\begin{tabular}{@{}p{4.6cm}|p{4.6cm}@{}}
\hline
\textbf{TEM o som} & \textbf{NÃO tem o som} \\
\hline
\textbf{bola} & \textbf{chá} \\
\textbf{ovo} & \textbf{pá} \\
\textbf{horta} & \textbf{mês} \\
\hline
\end{tabular}
\end{center}

\textbf{C) COMPLETE} --- Complete com a grafia do som e confira entre parênteses.

\begin{center}
\large p\underline{\hspace{0.7cm}} (pó) \quad n\underline{\hspace{0.7cm}} (nó) \quad s\underline{\hspace{0.7cm}} (só) \quad av\underline{\hspace{0.7cm}} (avó)
\end{center}

\textbf{D) OUÇA E COMPARE} --- Ó aberto e Ô fechado.

\begin{center}
\begin{tabular}{@{}p{4.2cm}p{6.2cm}p{1.6cm}@{}}
\hline
\textbf{Palavra} & \textbf{Compare} & \textbf{Marque} \\
\hline
\textbf{avó} & abre no fim [o aberto] & \textbf{SIM} \\
\textbf{avô} & fecha no fim [o] & \textbf{NÃO} \\
\textbf{pó} & abre [o aberto] & \textbf{SIM} \\
\hline
\end{tabular}
\end{center}

\caixapausa{Se precisar, pare aqui. Respire fundo e continue quando quiser.}
\caixapista{Lábios arredondados e bem abertos. Mão na frente da boca: ar em círculo.}
\end{exercicio}

\plancheta{Atividade de Som 7/38 --- /o/ --- AVÔ, BOLO e FOGO}
\begin{exercicio}[Som /o/ [[o]] --- Atividade de Som (não de letra nem de sílaba)]
\noindent\textbf{Nome:} \underline{\hspace{3.2cm}} \quad \textbf{Data:} \underline{\hspace{1.6cm}}

\noindent\textbf{\faMusic~O som /o/ {[o]}:} avô, bolo, fogo, coro, gosto.
\newline\textbf{\faPencilRuler~Quem faz esse som:} Ô fechado. \hfill \textbf{\faCheckCircle~Voz:} com voz.

\vspace{0.4em}
\textbf{A) CAÇA AO SOM} --- Leia em voz alta e compare os sons. Marque a coluna certa para cada palavra.

\begin{center}
\begin{tabular}{@{}ll@{}}
\fbox{\textbf{SIM}} \textit{avô} \quad \fbox{\textbf{NÃO}} \textit{pó} \\
\fbox{\textbf{SIM}} \textit{bolo} \quad \fbox{\textbf{NÃO}} \textit{pé} \\
\fbox{\textbf{SIM}} \textit{fogo} \quad \fbox{\textbf{NÃO}} \textit{pá} \\
\fbox{\textbf{SIM}} \textit{coro} \quad \fbox{\textbf{NÃO}} \textit{nó} \\
\end{tabular}
\end{center}

\textbf{B) SEPARE} --- Escreva as palavras na coluna certa.

\begin{center}
\begin{tabular}{@{}p{4.6cm}|p{4.6cm}@{}}
\hline
\textbf{TEM o som} & \textbf{NÃO tem o som} \\
\hline
\textbf{gosto} & \textbf{chá} \\
\textbf{avô} & \textbf{bola} \\
\textbf{coro} & \textbf{só} \\
\hline
\end{tabular}
\end{center}

\textbf{C) COMPLETE} --- Complete com a grafia do som e confira entre parênteses.

\begin{center}
\large av\underline{\hspace{0.7cm}} (avô) \quad b\underline{\hspace{0.7cm}}l\underline{\hspace{0.7cm}} (bolo) \quad f\underline{\hspace{0.7cm}}g\underline{\hspace{0.7cm}} (fogo) \quad c\underline{\hspace{0.7cm}}r\underline{\hspace{0.7cm}} (coro)
\end{center}

\textbf{D) OUÇA E COMPARE} --- Ô fechado e Ó aberto.

\begin{center}
\begin{tabular}{@{}p{4.2cm}p{6.2cm}p{1.6cm}@{}}
\hline
\textbf{Palavra} & \textbf{Compare} & \textbf{Marque} \\
\hline
\textbf{avô} & fecha [o] & \textbf{NÃO} \\
\textbf{avó} & abre [o aberto] & \textbf{SIM} \\
\textbf{bolo} & fecha [o] & \textbf{NÃO} \\
\hline
\end{tabular}
\end{center}

\caixapausa{Se precisar, pare aqui. Respire fundo e continue quando quiser.}
\caixapista{Lábios arredondados, boca mais fechada. Mão na garganta: vibra.}
\end{exercicio}

\plancheta{Atividade de Som 8/38 --- /u/ --- UVA, LUA e UNHA}
\begin{exercicio}[Som /u/ [[u]] --- Atividade de Som (não de letra nem de sílaba)]
\noindent\textbf{Nome:} \underline{\hspace{3.2cm}} \quad \textbf{Data:} \underline{\hspace{1.6cm}}

\noindent\textbf{\faMusic~O som /u/ {[u]}:} uva, lua, unha, urubu, chuchu.
\newline\textbf{\faPencilRuler~Quem faz esse som:} U. \hfill \textbf{\faCheckCircle~Voz:} com voz.

\vspace{0.4em}
\textbf{A) CAÇA AO SOM} --- Leia em voz alta e compare os sons. Marque a coluna certa para cada palavra.

\begin{center}
\begin{tabular}{@{}ll@{}}
\fbox{\textbf{SIM}} \textit{uva} \quad \fbox{\textbf{NÃO}} \textit{pé} \\
\fbox{\textbf{SIM}} \textit{lua} \quad \fbox{\textbf{NÃO}} \textit{mês} \\
\fbox{\textbf{SIM}} \textit{unha} \quad \fbox{\textbf{NÃO}} \textit{pá} \\
\fbox{\textbf{SIM}} \textit{urubu} \quad \fbox{\textbf{NÃO}} \textit{café} \\
\end{tabular}
\end{center}

\textbf{B) SEPARE} --- Escreva as palavras na coluna certa.

\begin{center}
\begin{tabular}{@{}p{4.6cm}|p{4.6cm}@{}}
\hline
\textbf{TEM o som} & \textbf{NÃO tem o som} \\
\hline
\textbf{chuchu} & \textbf{avó} \\
\textbf{jacaré} & \textbf{dedo} \\
\textbf{futuro} & \textbf{sopa} \\
\hline
\end{tabular}
\end{center}

\textbf{C) COMPLETE} --- Complete com a grafia do som e confira entre parênteses.

\begin{center}
\large \underline{\hspace{0.7cm}}va (uva) \quad l\underline{\hspace{0.7cm}}a (lua) \quad \underline{\hspace{0.7cm}}nha (unha)
\end{center}

\textbf{D) OUÇA E COMPARE} --- U forte e U rápido.

\begin{center}
\begin{tabular}{@{}p{4.2cm}p{6.2cm}p{1.6cm}@{}}
\hline
\textbf{Palavra} & \textbf{Compare} & \textbf{Marque} \\
\hline
\textbf{lua} & U forte no meio & \textbf{SIM} \\
\textbf{pau} & U rápido no fim [w] & \textbf{NÃO} \\
\textbf{rua} & U forte & \textbf{SIM} \\
\hline
\end{tabular}
\end{center}

\caixapausa{Se precisar, pare aqui. Respire fundo e continue quando quiser.}
\caixapista{Lábios em bico, som grave e fechado. Mão na frente da boca: ar focado.}
\end{exercicio}

\plancheta{Atividade de Som 9/38 --- /ã/ --- ANÃ, MAÇÃ e ROMÃ}
\begin{exercicio}[Som /ã/ [[a~]] --- Atividade de Som (não de letra nem de sílaba)]
\noindent\textbf{Nome:} \underline{\hspace{3.2cm}} \quad \textbf{Data:} \underline{\hspace{1.6cm}}

\noindent\textbf{\faMusic~O som /ã/ {[a~]}:} anã, maçã, romã, ímã.
\newline\textbf{\faPencilRuler~Quem faz esse som:} Ã (til). \hfill \textbf{\faCheckCircle~Voz:} com voz, nasal.

\vspace{0.4em}
\textbf{A) CAÇA AO SOM} --- Leia cada par com o dedo no nariz: marca NARIZ para quem faz o nariz vibrar e BOCA para quem não faz.

\begin{center}
\begin{tabular}{@{}ll@{}}
\fbox{\textbf{SIM}} \textit{anã} \quad \fbox{\textbf{NÃO}} \textit{bola} \\
\fbox{\textbf{SIM}} \textit{maçã} \quad \fbox{\textbf{NÃO}} \textit{pé} \\
\fbox{\textbf{SIM}} \textit{romã} \quad \fbox{\textbf{NÃO}} \textit{café} \\
\fbox{\textbf{SIM}} \textit{ímã} \quad \fbox{\textbf{NÃO}} \textit{chá} \\
\end{tabular}
\end{center}

\textbf{B) SEPARE} --- Escreva as palavras na coluna certa.

\begin{center}
\begin{tabular}{@{}p{4.6cm}|p{4.6cm}@{}}
\hline
\textbf{TEM o som} & \textbf{NÃO tem o som} \\
\hline
\textbf{maçã} & \textbf{lá} \\
\textbf{anã} & \textbf{pó} \\
\textbf{romã} & \textbf{só} \\
\hline
\end{tabular}
\end{center}

\textbf{C) COMPLETE} --- Complete com a grafia do som e confira entre parênteses.

\begin{center}
\large an\underline{\hspace{0.7cm}} (anã) \quad maç\underline{\hspace{0.7cm}} (maçã) \quad rom\underline{\hspace{0.7cm}} (romã) \quad ím\underline{\hspace{0.7cm}} (ímã)
\end{center}

\textbf{D) OUÇA E COMPARE} --- Oral e nasal.

\begin{center}
\begin{tabular}{@{}p{4.2cm}p{6.2cm}p{1.6cm}@{}}
\hline
\textbf{Palavra} & \textbf{Compare} & \textbf{Marque} \\
\hline
\textbf{lá} & nariz não vibra & \textbf{NÃO} \\
\textbf{lã} & nariz vibra [ã] & \textbf{SIM} \\

\hline
\end{tabular}
\end{center}

\caixapausa{Se precisar, pare aqui. Respire fundo e continue quando quiser.}
\caixapista{Boca aberta e ar saindo pelo nariz. Mão no nariz: sente a vibração.}
\end{exercicio}

\plancheta{Atividade de Som 10/38 --- /ã/ fraco --- CANTA, TAMPA e BANDA}
\begin{exercicio}[Som /ã/ fraco [[6~]] --- Atividade de Som (não de letra nem de sílaba)]
\noindent\textbf{Nome:} \underline{\hspace{3.2cm}} \quad \textbf{Data:} \underline{\hspace{1.6cm}}

\noindent\textbf{\faMusic~O som /ã/ fraco {[6~]}:} canta, tampa, banda, campo, andar.
\newline\textbf{\faPencilRuler~Quem faz esse som:} AN / AM (fraco). \hfill \textbf{\faCheckCircle~Voz:} com voz, nasal.

\vspace{0.4em}
\textbf{A) CAÇA AO SOM} --- Leia cada par com o dedo no nariz: marca NARIZ para quem faz o nariz vibrar e BOCA para quem não faz.

\begin{center}
\begin{tabular}{@{}ll@{}}
\fbox{\textbf{SIM}} \textit{canta} \quad \fbox{\textbf{NÃO}} \textit{pato} \\
\fbox{\textbf{SIM}} \textit{tampa} \quad \fbox{\textbf{NÃO}} \textit{bola} \\
\fbox{\textbf{SIM}} \textit{banda} \quad \fbox{\textbf{NÃO}} \textit{pé} \\
\fbox{\textbf{SIM}} \textit{campo} \quad \fbox{\textbf{NÃO}} \textit{café} \\
\end{tabular}
\end{center}

\textbf{B) SEPARE} --- Escreva as palavras na coluna certa.

\begin{center}
\begin{tabular}{@{}p{4.6cm}|p{4.6cm}@{}}
\hline
\textbf{TEM o som} & \textbf{NÃO tem o som} \\
\hline
\textbf{santo} & \textbf{sopa} \\
\textbf{tampa} & \textbf{dedo} \\
\textbf{andar} & \textbf{mala} \\
\hline
\end{tabular}
\end{center}

\textbf{C) COMPLETE} --- Complete com a grafia do som e confira entre parênteses.

\begin{center}
\large c\underline{\hspace{1.1cm}}ta (canta) \quad t\underline{\hspace{1.1cm}}pa (tampa) \quad b\underline{\hspace{1.1cm}}da (banda) \quad c\underline{\hspace{1.1cm}}po (campo)
\end{center}

\textbf{D) OUÇA E COMPARE} --- Oral e nasal.

\begin{center}
\begin{tabular}{@{}p{4.2cm}p{6.2cm}p{1.6cm}@{}}
\hline
\textbf{Palavra} & \textbf{Compare} & \textbf{Marque} \\
\hline
\textbf{tapa} & sem nasal [a] & \textbf{NÃO} \\
\textbf{tampa} & com nasal [ã fraco] & \textbf{SIM} \\
\textbf{pato} & sem nasal & \textbf{NÃO} \\
\hline
\end{tabular}
\end{center}

\caixapausa{Se precisar, pare aqui. Respire fundo e continue quando quiser.}
\caixapista{Boca mais fechada, ar pelo nariz. Mão no nariz: vibra forte.}
\end{exercicio}

\plancheta{Atividade de Som 11/38 --- /e/ nasal --- TEMPO, SEMPRE e CENTO}
\begin{exercicio}[Som /e/ nasal [[e~]] --- Atividade de Som (não de letra nem de sílaba)]
\noindent\textbf{Nome:} \underline{\hspace{3.2cm}} \quad \textbf{Data:} \underline{\hspace{1.6cm}}

\noindent\textbf{\faMusic~O som /e/ nasal {[e~]}:} tempo, sempre, cento, pente, cem.
\newline\textbf{\faPencilRuler~Quem faz esse som:} EN / EM. \hfill \textbf{\faCheckCircle~Voz:} com voz, nasal.

\vspace{0.4em}
\textbf{A) CAÇA AO SOM} --- Leia cada par com o dedo no nariz: marca NARIZ para quem faz o nariz vibrar e BOCA para quem não faz.

\begin{center}
\begin{tabular}{@{}ll@{}}
\fbox{\textbf{SIM}} \textit{tempo} \quad \fbox{\textbf{NÃO}} \textit{dedo} \\
\fbox{\textbf{SIM}} \textit{sempre} \quad \fbox{\textbf{NÃO}} \textit{pé} \\
\fbox{\textbf{SIM}} \textit{cento} \quad \fbox{\textbf{NÃO}} \textit{café} \\
\fbox{\textbf{SIM}} \textit{pente} \quad \fbox{\textbf{NÃO}} \textit{bola} \\
\end{tabular}
\end{center}

\textbf{B) SEPARE} --- Escreva as palavras na coluna certa.

\begin{center}
\begin{tabular}{@{}p{4.6cm}|p{4.6cm}@{}}
\hline
\textbf{TEM o som} & \textbf{NÃO tem o som} \\
\hline
\textbf{cem} & \textbf{chá} \\
\textbf{tempo} & \textbf{lá} \\
\textbf{pente} & \textbf{pó} \\
\hline
\end{tabular}
\end{center}

\textbf{C) COMPLETE} --- Complete com a grafia do som e confira entre parênteses.

\begin{center}
\large t\underline{\hspace{1.1cm}}po (tempo) \quad c\underline{\hspace{1.1cm}}to (cento) \quad p\underline{\hspace{1.1cm}}te (pente)
\end{center}

\textbf{D) OUÇA E COMPARE} --- Oral e nasal.

\begin{center}
\begin{tabular}{@{}p{4.2cm}p{6.2cm}p{1.6cm}@{}}
\hline
\textbf{Palavra} & \textbf{Compare} & \textbf{Marque} \\
\hline
\textbf{pé} & sem nasal [e aberto] & \textbf{NÃO} \\
\textbf{pente} & com nasal [e nasal] & \textbf{SIM} \\
\textbf{chá} & sem nasal & \textbf{NÃO} \\
\hline
\end{tabular}
\end{center}

\caixapausa{Se precisar, pare aqui. Respire fundo e continue quando quiser.}
\caixapista{Som de E + ar pelo nariz. Mão no nariz: vibra.}
\end{exercicio}

\plancheta{Atividade de Som 12/38 --- /i/ nasal --- SIM, FIM e LIMPO}
\begin{exercicio}[Som /i/ nasal [[i~]] --- Atividade de Som (não de letra nem de sílaba)]
\noindent\textbf{Nome:} \underline{\hspace{3.2cm}} \quad \textbf{Data:} \underline{\hspace{1.6cm}}

\noindent\textbf{\faMusic~O som /i/ nasal {[i~]}:} sim, fim, mim, limpo, pinto.
\newline\textbf{\faPencilRuler~Quem faz esse som:} IN / IM. \hfill \textbf{\faCheckCircle~Voz:} com voz, nasal.

\vspace{0.4em}
\textbf{A) CAÇA AO SOM} --- Leia cada par com o dedo no nariz: marca NARIZ para quem faz o nariz vibrar e BOCA para quem não faz.

\begin{center}
\begin{tabular}{@{}ll@{}}
\fbox{\textbf{SIM}} \textit{sim} \quad \fbox{\textbf{NÃO}} \textit{si} \\
\fbox{\textbf{SIM}} \textit{fim} \quad \fbox{\textbf{NÃO}} \textit{lá} \\
\fbox{\textbf{SIM}} \textit{mim} \quad \fbox{\textbf{NÃO}} \textit{pá} \\
\fbox{\textbf{SIM}} \textit{limpo} \quad \fbox{\textbf{NÃO}} \textit{sol} \\
\end{tabular}
\end{center}

\textbf{B) SEPARE} --- Escreva as palavras na coluna certa.

\begin{center}
\begin{tabular}{@{}p{4.6cm}|p{4.6cm}@{}}
\hline
\textbf{TEM o som} & \textbf{NÃO tem o som} \\
\hline
\textbf{pinto} & \textbf{chá} \\
\textbf{sim} & \textbf{pó} \\
\textbf{fim} & \textbf{só} \\
\hline
\end{tabular}
\end{center}

\textbf{C) COMPLETE} --- Complete com a grafia do som e confira entre parênteses.

\begin{center}
\large s\underline{\hspace{0.7cm}}m (sim) \quad f\underline{\hspace{0.7cm}}m (fim) \quad m\underline{\hspace{0.7cm}}m (mim) \quad l\underline{\hspace{1.1cm}}po (limpo)
\end{center}

\textbf{D) OUÇA E COMPARE} --- Oral e nasal.

\begin{center}
\begin{tabular}{@{}p{4.2cm}p{6.2cm}p{1.6cm}@{}}
\hline
\textbf{Palavra} & \textbf{Compare} & \textbf{Marque} \\
\hline
\textbf{si} & sem nasal [i] & \textbf{NÃO} \\
\textbf{sim} & com nasal [i nasal] & \textbf{SIM} \\
\textbf{vi} & sem nasal & \textbf{NÃO} \\
\hline
\end{tabular}
\end{center}

\caixapausa{Se precisar, pare aqui. Respire fundo e continue quando quiser.}
\caixapista{Som de I + ar pelo nariz. Mão no nariz: vibra.}
\end{exercicio}

\plancheta{Atividade de Som 13/38 --- /o/ nasal --- BOM, SOM e PONTE}
\begin{exercicio}[Som /o/ nasal [[o~]] --- Atividade de Som (não de letra nem de sílaba)]
\noindent\textbf{Nome:} \underline{\hspace{3.2cm}} \quad \textbf{Data:} \underline{\hspace{1.6cm}}

\noindent\textbf{\faMusic~O som /o/ nasal {[o~]}:} bom, som, tom, ponte, pomba.
\newline\textbf{\faPencilRuler~Quem faz esse som:} ON / OM. \hfill \textbf{\faCheckCircle~Voz:} com voz, nasal.

\vspace{0.4em}
\textbf{A) CAÇA AO SOM} --- Leia cada par com o dedo no nariz: marca NARIZ para quem faz o nariz vibrar e BOCA para quem não faz.

\begin{center}
\begin{tabular}{@{}ll@{}}
\fbox{\textbf{SIM}} \textit{bom} \quad \fbox{\textbf{NÃO}} \textit{pá} \\
\fbox{\textbf{SIM}} \textit{som} \quad \fbox{\textbf{NÃO}} \textit{pé} \\
\fbox{\textbf{SIM}} \textit{tom} \quad \fbox{\textbf{NÃO}} \textit{café} \\
\fbox{\textbf{SIM}} \textit{ponte} \quad \fbox{\textbf{NÃO}} \textit{sol} \\
\end{tabular}
\end{center}

\textbf{B) SEPARE} --- Escreva as palavras na coluna certa.

\begin{center}
\begin{tabular}{@{}p{4.6cm}|p{4.6cm}@{}}
\hline
\textbf{TEM o som} & \textbf{NÃO tem o som} \\
\hline
\textbf{pomba} & \textbf{chá} \\
\textbf{bom} & \textbf{pó} \\
\textbf{tom} & \textbf{só} \\
\hline
\end{tabular}
\end{center}

\textbf{C) COMPLETE} --- Complete com a grafia do som e confira entre parênteses.

\begin{center}
\large b\underline{\hspace{0.7cm}}m (bom) \quad s\underline{\hspace{0.7cm}}m (som) \quad t\underline{\hspace{0.7cm}}m (tom) \quad p\underline{\hspace{1.1cm}}te (ponte)
\end{center}

\textbf{D) OUÇA E COMPARE} --- Oral e nasal.

\begin{center}
\begin{tabular}{@{}p{4.2cm}p{6.2cm}p{1.6cm}@{}}
\hline
\textbf{Palavra} & \textbf{Compare} & \textbf{Marque} \\
\hline
\textbf{só} & sem nasal [o aberto] & \textbf{NÃO} \\
\textbf{som} & com nasal [õ] & \textbf{SIM} \\
\textbf{pó} & sem nasal & \textbf{NÃO} \\
\hline
\end{tabular}
\end{center}

\caixapausa{Se precisar, pare aqui. Respire fundo e continue quando quiser.}
\caixapista{Som de O + ar pelo nariz. Mão no nariz: vibra.}
\end{exercicio}

\plancheta{Atividade de Som 14/38 --- /u/ nasal --- UM, ALGUM e JEJUM}
\begin{exercicio}[Som /u/ nasal [[u~]] --- Atividade de Som (não de letra nem de sílaba)]
\noindent\textbf{Nome:} \underline{\hspace{3.2cm}} \quad \textbf{Data:} \underline{\hspace{1.6cm}}

\noindent\textbf{\faMusic~O som /u/ nasal {[u~]}:} um, algum, jejum, comum, atum.
\newline\textbf{\faPencilRuler~Quem faz esse som:} UN / UM. \hfill \textbf{\faCheckCircle~Voz:} com voz, nasal.

\vspace{0.4em}
\textbf{A) CAÇA AO SOM} --- Leia cada par com o dedo no nariz: marca NARIZ para quem faz o nariz vibrar e BOCA para quem não faz.

\begin{center}
\begin{tabular}{@{}ll@{}}
\fbox{\textbf{SIM}} \textit{um} \quad \fbox{\textbf{NÃO}} \textit{pé} \\
\fbox{\textbf{SIM}} \textit{algum} \quad \fbox{\textbf{NÃO}} \textit{café} \\
\fbox{\textbf{SIM}} \textit{jejum} \quad \fbox{\textbf{NÃO}} \textit{pá} \\
\fbox{\textbf{SIM}} \textit{comum} \quad \fbox{\textbf{NÃO}} \textit{sol} \\
\end{tabular}
\end{center}

\textbf{B) SEPARE} --- Escreva as palavras na coluna certa.

\begin{center}
\begin{tabular}{@{}p{4.6cm}|p{4.6cm}@{}}
\hline
\textbf{TEM o som} & \textbf{NÃO tem o som} \\
\hline
\textbf{atum} & \textbf{chá} \\
\textbf{um} & \textbf{pó} \\
\textbf{jejum} & \textbf{só} \\
\hline
\end{tabular}
\end{center}

\textbf{C) COMPLETE} --- Complete com a grafia do som e confira entre parênteses.

\begin{center}
\large \underline{\hspace{0.7cm}}m (um) \quad alg\underline{\hspace{0.7cm}}m (algum) \quad jej\underline{\hspace{0.7cm}}m (jejum) \quad com\underline{\hspace{0.7cm}}m (comum)
\end{center}

\textbf{D) OUÇA E COMPARE} --- Oral e nasal.

\begin{center}
\begin{tabular}{@{}p{4.2cm}p{6.2cm}p{1.6cm}@{}}
\hline
\textbf{Palavra} & \textbf{Compare} & \textbf{Marque} \\
\hline
\textbf{u} & sem nasal [u] & \textbf{NÃO} \\
\textbf{um} & com nasal [u nasal] & \textbf{SIM} \\
\textbf{tu} & sem nasal & \textbf{NÃO} \\
\hline
\end{tabular}
\end{center}

\caixapausa{Se precisar, pare aqui. Respire fundo e continue quando quiser.}
\caixapista{Som de U + ar pelo nariz. Mão no nariz: vibra.}
\end{exercicio}

\plancheta{Atividade de Som 15/38 --- /j/ --- PAI, VAI e BOI}
\begin{exercicio}[Som /j/ [[j]] --- Atividade de Som (não de letra nem de sílaba)]
\noindent\textbf{Nome:} \underline{\hspace{3.2cm}} \quad \textbf{Data:} \underline{\hspace{1.6cm}}

\noindent\textbf{\faMusic~O som /j/ {[j]}:} pai, vai, cai, boi, papai, herói.
\newline\textbf{\faPencilRuler~Quem faz esse som:} I e E semivogais. \hfill \textbf{\faCheckCircle~Voz:} com voz, rápido.

\vspace{0.4em}
\textbf{A) CAÇA AO SOM} --- Leia em voz alta e compare os sons. Marque a coluna certa para cada palavra.

\begin{center}
\begin{tabular}{@{}ll@{}}
\fbox{\textbf{SIM}} \textit{pai} \quad \fbox{\textbf{NÃO}} \textit{pá} \\
\fbox{\textbf{SIM}} \textit{vai} \quad \fbox{\textbf{NÃO}} \textit{pé} \\
\fbox{\textbf{SIM}} \textit{cai} \quad \fbox{\textbf{NÃO}} \textit{café} \\
\fbox{\textbf{SIM}} \textit{boi} \quad \fbox{\textbf{NÃO}} \textit{chá} \\
\end{tabular}
\end{center}

\textbf{B) SEPARE} --- Escreva as palavras na coluna certa.

\begin{center}
\begin{tabular}{@{}p{4.6cm}|p{4.6cm}@{}}
\hline
\textbf{TEM o som} & \textbf{NÃO tem o som} \\
\hline
\textbf{papai} & \textbf{sol} \\
\textbf{herói} & \textbf{lá} \\
\textbf{mãe} & \textbf{pó} \\
\hline
\end{tabular}
\end{center}

\textbf{C) COMPLETE} --- Complete com a grafia do som e confira entre parênteses.

\begin{center}
\large pa\underline{\hspace{0.7cm}} (pai) \quad va\underline{\hspace{0.7cm}} (vai) \quad ca\underline{\hspace{0.7cm}} (cai) \quad bo\underline{\hspace{0.7cm}} (boi)
\end{center}

\textbf{D) OUÇA E COMPARE} --- I forte e I rápido.

\begin{center}
\begin{tabular}{@{}p{4.2cm}p{6.2cm}p{1.6cm}@{}}
\hline
\textbf{Palavra} & \textbf{Compare} & \textbf{Marque} \\
\hline
\textbf{pipa} & I forte no meio [i] & \textbf{NÃO} \\
\textbf{pai} & I rápido no fim [j] & \textbf{SIM} \\
\textbf{sino} & I forte & \textbf{NÃO} \\
\hline
\end{tabular}
\end{center}

\caixapausa{Se precisar, pare aqui. Respire fundo e continue quando quiser.}
\caixapista{Lábios em sorriso; som rápido no fim da sílaba. Mão na garganta: vibra.}
\end{exercicio}

\plancheta{Atividade de Som 16/38 --- /w/ --- PAU, CÉU e SOL}
\begin{exercicio}[Som /w/ [[w]] --- Atividade de Som (não de letra nem de sílaba)]
\noindent\textbf{Nome:} \underline{\hspace{3.2cm}} \quad \textbf{Data:} \underline{\hspace{1.6cm}}

\noindent\textbf{\faMusic~O som /w/ {[w]}:} pau, mau, céu, réu; sol, mel, farol; qual, quando.
\newline\textbf{\faPencilRuler~Quem faz esse som:} U e L semivogais. \hfill \textbf{\faCheckCircle~Voz:} com voz, rápido.

\vspace{0.4em}
\textbf{A) CAÇA AO SOM} --- Leia em voz alta e compare os sons. Marque a coluna certa para cada palavra.

\begin{center}
\begin{tabular}{@{}ll@{}}
\fbox{\textbf{SIM}} \textit{pau} \quad \fbox{\textbf{NÃO}} \textit{pá} \\
\fbox{\textbf{SIM}} \textit{mau} \quad \fbox{\textbf{NÃO}} \textit{pé} \\
\fbox{\textbf{SIM}} \textit{céu} \quad \fbox{\textbf{NÃO}} \textit{café} \\
\fbox{\textbf{SIM}} \textit{sol} \quad \fbox{\textbf{NÃO}} \textit{chá} \\
\end{tabular}
\end{center}

\textbf{B) SEPARE} --- Escreva as palavras na coluna certa.

\begin{center}
\begin{tabular}{@{}p{4.6cm}|p{4.6cm}@{}}
\hline
\textbf{TEM o som} & \textbf{NÃO tem o som} \\
\hline
\textbf{mel} & \textbf{lá} \\
\textbf{réu} & \textbf{pó} \\
\textbf{farol} & \textbf{só} \\
\hline
\end{tabular}
\end{center}

\textbf{C) COMPLETE} --- Complete com a grafia do som e confira entre parênteses.

\begin{center}
\large pa\underline{\hspace{0.7cm}} (pau) \quad ma\underline{\hspace{0.7cm}} (mau) \quad cé\underline{\hspace{0.7cm}} (céu) \quad se\underline{\hspace{0.7cm}}l (seu)
\end{center}

\textbf{D) OUÇA E COMPARE} --- U forte e U rápido.

\begin{center}
\begin{tabular}{@{}p{4.2cm}p{6.2cm}p{1.6cm}@{}}
\hline
\textbf{Palavra} & \textbf{Compare} & \textbf{Marque} \\
\hline
\textbf{lua} & U forte [u] & \textbf{NÃO} \\
\textbf{pau} & U rápido [w] & \textbf{SIM} \\
\textbf{rua} & U forte & \textbf{NÃO} \\
\hline
\end{tabular}
\end{center}

\caixapausa{Se precisar, pare aqui. Respire fundo e continue quando quiser.}
\caixapista{Lábios em bico; som rápido. Mão na frente da boca: ar rápido.}
\end{exercicio}

\plancheta{Atividade de Som 17/38 --- /p/ --- PATO, PIPA e POTE}
\begin{exercicio}[Som /p/ [[p]] --- Atividade de Som (não de letra nem de sílaba)]
\noindent\textbf{Nome:} \underline{\hspace{3.2cm}} \quad \textbf{Data:} \underline{\hspace{1.6cm}}

\noindent\textbf{\faMusic~O som /p/ {[p]}:} pato, papai, pipa, pote, pulo, sapo.
\newline\textbf{\faPencilRuler~Quem faz esse som:} P. \hfill \textbf{\faCheckCircle~Voz:} sem voz.

\vspace{0.4em}
\textbf{A) CAÇA AO SOM} --- Leia cada par com a mão na garganta: marca VIBRA para quem vibra e NÃO VIBRA para quem não vibra.

\begin{center}
\begin{tabular}{@{}ll@{}}
\fbox{\textbf{SIM}} \textit{pato} \quad \fbox{\textbf{NÃO}} \textit{bolo} \\
\fbox{\textbf{SIM}} \textit{papai} \quad \fbox{\textbf{NÃO}} \textit{dedo} \\
\fbox{\textbf{SIM}} \textit{pipa} \quad \fbox{\textbf{NÃO}} \textit{casa} \\
\fbox{\textbf{SIM}} \textit{pote} \quad \fbox{\textbf{NÃO}} \textit{mala} \\
\end{tabular}
\end{center}

\textbf{B) SEPARE} --- Escreva as palavras na coluna certa.

\begin{center}
\begin{tabular}{@{}p{4.6cm}|p{4.6cm}@{}}
\hline
\textbf{TEM o som} & \textbf{NÃO tem o som} \\
\hline
\textbf{sapo} & \textbf{gato} \\
\textbf{pulo} & \textbf{teto} \\
\textbf{pipa} & \textbf{bola} \\
\hline
\end{tabular}
\end{center}

\textbf{C) COMPLETE} --- Complete com a grafia do som e confira entre parênteses.

\begin{center}
\large \underline{\hspace{0.7cm}}ato (pato) \quad \underline{\hspace{0.7cm}}ipa (pipa) \quad \underline{\hspace{0.7cm}}ote (pote) \quad \underline{\hspace{0.7cm}}ulo (pulo)
\end{center}

\textbf{D) OUÇA E COMPARE} --- P sem voz e B com voz.

\begin{center}
\begin{tabular}{@{}p{4.2cm}p{6.2cm}p{1.6cm}@{}}
\hline
\textbf{Palavra} & \textbf{Compare} & \textbf{Marque} \\
\hline
\textbf{pato} & sem voz & \textbf{NÃO} \\
\textbf{bato} & com voz & \textbf{SIM} \\
\textbf{pote} & sem voz & \textbf{NÃO} \\
\hline
\end{tabular}
\end{center}

\caixapausa{Se precisar, pare aqui. Respire fundo e continue quando quiser.}
\caixapista{Lábios fechados que explodem de repente. Mão na garganta: não vibra (som surdo).}
\end{exercicio}

\plancheta{Atividade de Som 18/38 --- /b/ --- BOLA, BABA e BEBÊ}
\begin{exercicio}[Som /b/ [[b]] --- Atividade de Som (não de letra nem de sílaba)]
\noindent\textbf{Nome:} \underline{\hspace{3.2cm}} \quad \textbf{Data:} \underline{\hspace{1.6cm}}

\noindent\textbf{\faMusic~O som /b/ {[b]}:} bola, baba, bebê, bota, botão.
\newline\textbf{\faPencilRuler~Quem faz esse som:} B. \hfill \textbf{\faCheckCircle~Voz:} com voz.

\vspace{0.4em}
\textbf{A) CAÇA AO SOM} --- Leia cada par com a mão na garganta: marca VIBRA para quem vibra e NÃO VIBRA para quem não vibra.

\begin{center}
\begin{tabular}{@{}ll@{}}
\fbox{\textbf{SIM}} \textit{bola} \quad \fbox{\textbf{NÃO}} \textit{pato} \\
\fbox{\textbf{SIM}} \textit{baba} \quad \fbox{\textbf{NÃO}} \textit{copo} \\
\fbox{\textbf{SIM}} \textit{bebê} \quad \fbox{\textbf{NÃO}} \textit{mala} \\
\fbox{\textbf{SIM}} \textit{bota} \quad \fbox{\textbf{NÃO}} \textit{dedo} \\
\end{tabular}
\end{center}

\textbf{B) SEPARE} --- Escreva as palavras na coluna certa.

\begin{center}
\begin{tabular}{@{}p{4.6cm}|p{4.6cm}@{}}
\hline
\textbf{TEM o som} & \textbf{NÃO tem o som} \\
\hline
\textbf{banana} & \textbf{sapo} \\
\textbf{bola} & \textbf{teto} \\
\textbf{bebê} & \textbf{pipa} \\
\hline
\end{tabular}
\end{center}

\textbf{C) COMPLETE} --- Complete com a grafia do som e confira entre parênteses.

\begin{center}
\large \underline{\hspace{0.7cm}}ola (bola) \quad \underline{\hspace{0.7cm}}aba (baba) \quad \underline{\hspace{0.7cm}}ebê (bebê) \quad \underline{\hspace{0.7cm}}ota (bota)
\end{center}

\textbf{D) OUÇA E COMPARE} --- B com voz e P sem voz.

\begin{center}
\begin{tabular}{@{}p{4.2cm}p{6.2cm}p{1.6cm}@{}}
\hline
\textbf{Palavra} & \textbf{Compare} & \textbf{Marque} \\
\hline
\textbf{bato} & com voz & \textbf{SIM} \\
\textbf{pato} & sem voz & \textbf{NÃO} \\
\textbf{bote} & com voz & \textbf{SIM} \\
\hline
\end{tabular}
\end{center}

\caixapausa{Se precisar, pare aqui. Respire fundo e continue quando quiser.}
\caixapista{Lábios fechados que explodem COM voz. Mão na garganta: vibra (som sonoro).}
\end{exercicio}

\plancheta{Atividade de Som 19/38 --- /t/ --- TATU, TETO e TUDO}
\begin{exercicio}[Som /t/ [[t]] --- Atividade de Som (não de letra nem de sílaba)]
\noindent\textbf{Nome:} \underline{\hspace{3.2cm}} \quad \textbf{Data:} \underline{\hspace{1.6cm}}

\noindent\textbf{\faMusic~O som /t/ {[t]}:} tatu, teto, tudo, tela, pato, pote.
\newline\textbf{\faPencilRuler~Quem faz esse som:} T. \hfill \textbf{\faCheckCircle~Voz:} sem voz (varia em TIA, LEITE — veja nota).

\vspace{0.4em}
\textbf{A) CAÇA AO SOM} --- Leia cada par com a mão na garganta: marca VIBRA para quem vibra e NÃO VIBRA para quem não vibra.

\begin{center}
\begin{tabular}{@{}ll@{}}
\fbox{\textbf{SIM}} \textit{tatu} \quad \fbox{\textbf{NÃO}} \textit{dedo} \\
\fbox{\textbf{SIM}} \textit{teto} \quad \fbox{\textbf{NÃO}} \textit{bola} \\
\fbox{\textbf{SIM}} \textit{tudo} \quad \fbox{\textbf{NÃO}} \textit{pé} \\
\fbox{\textbf{SIM}} \textit{tela} \quad \fbox{\textbf{NÃO}} \textit{casa} \\
\end{tabular}
\end{center}

\textbf{B) SEPARE} --- Escreva as palavras na coluna certa.

\begin{center}
\begin{tabular}{@{}p{4.6cm}|p{4.6cm}@{}}
\hline
\textbf{TEM o som} & \textbf{NÃO tem o som} \\
\hline
\textbf{pato} & \textbf{gato} \\
\textbf{pote} & \textbf{mala} \\
\textbf{tela} & \textbf{sapo} \\
\hline
\end{tabular}
\end{center}

\textbf{C) COMPLETE} --- Complete com a grafia do som e confira entre parênteses.

\begin{center}
\large \underline{\hspace{0.7cm}}atu (tatu) \quad \underline{\hspace{0.7cm}}eto (teto) \quad \underline{\hspace{0.7cm}}udo (tudo) \quad \underline{\hspace{0.7cm}}ela (tela)
\end{center}

\textbf{D) OUÇA E COMPARE} --- T sem voz e D com voz.

\begin{center}
\begin{tabular}{@{}p{4.2cm}p{6.2cm}p{1.6cm}@{}}
\hline
\textbf{Palavra} & \textbf{Compare} & \textbf{Marque} \\
\hline
\textbf{tia} & sem voz [t] & \textbf{NÃO} \\
\textbf{dia} & com voz [d] & \textbf{SIM} \\
\textbf{teto} & sem voz & \textbf{NÃO} \\
\hline
\end{tabular}
\end{center}

\caixapausa{Se precisar, pare aqui. Respire fundo e continue quando quiser.}
\caixapista{Língua toca atrás dos dentes superiores. Mão na garganta: não vibra (som surdo).}
\end{exercicio}

\plancheta{Atividade de Som 20/38 --- /d/ --- DADO, DEDO e DOCA}
\begin{exercicio}[Som /d/ [[d]] --- Atividade de Som (não de letra nem de sílaba)]
\noindent\textbf{Nome:} \underline{\hspace{3.2cm}} \quad \textbf{Data:} \underline{\hspace{1.6cm}}

\noindent\textbf{\faMusic~O som /d/ {[d]}:} dado, dedo, doca, duna, dia, cidade.
\newline\textbf{\faPencilRuler~Quem faz esse som:} D. \hfill \textbf{\faCheckCircle~Voz:} com voz (varia em DIA, DICA — veja nota).

\vspace{0.4em}
\textbf{A) CAÇA AO SOM} --- Leia cada par com a mão na garganta: marca VIBRA para quem vibra e NÃO VIBRA para quem não vibra.

\begin{center}
\begin{tabular}{@{}ll@{}}
\fbox{\textbf{SIM}} \textit{dado} \quad \fbox{\textbf{NÃO}} \textit{pato} \\
\fbox{\textbf{SIM}} \textit{dedo} \quad \fbox{\textbf{NÃO}} \textit{bola} \\
\fbox{\textbf{SIM}} \textit{doca} \quad \fbox{\textbf{NÃO}} \textit{pé} \\
\fbox{\textbf{SIM}} \textit{duna} \quad \fbox{\textbf{NÃO}} \textit{casa} \\
\end{tabular}
\end{center}

\textbf{B) SEPARE} --- Escreva as palavras na coluna certa.

\begin{center}
\begin{tabular}{@{}p{4.6cm}|p{4.6cm}@{}}
\hline
\textbf{TEM o som} & \textbf{NÃO tem o som} \\
\hline
\textbf{cada} & \textbf{gato} \\
\textbf{dia} & \textbf{pipa} \\
\textbf{duna} & \textbf{sapo} \\
\hline
\end{tabular}
\end{center}

\textbf{C) COMPLETE} --- Complete com a grafia do som e confira entre parênteses.

\begin{center}
\large \underline{\hspace{0.7cm}}ado (dado) \quad \underline{\hspace{0.7cm}}edo (dedo) \quad \underline{\hspace{0.7cm}}oca (doca) \quad \underline{\hspace{0.7cm}}una (duna)
\end{center}

\textbf{D) OUÇA E COMPARE} --- D com voz e T sem voz.

\begin{center}
\begin{tabular}{@{}p{4.2cm}p{6.2cm}p{1.6cm}@{}}
\hline
\textbf{Palavra} & \textbf{Compare} & \textbf{Marque} \\
\hline
\textbf{dia} & com voz [d] & \textbf{SIM} \\
\textbf{tia} & sem voz [t] & \textbf{NÃO} \\
\textbf{dedo} & com voz & \textbf{SIM} \\
\hline
\end{tabular}
\end{center}

\caixapausa{Se precisar, pare aqui. Respire fundo e continue quando quiser.}
\caixapista{Língua toca atrás dos dentes superiores, COM voz. Mão na garganta: vibra (som sonoro).}
\end{exercicio}

\plancheta{Atividade de Som 21/38 --- /k/ --- CASA, QUEIJO e KART}
\begin{exercicio}[Som /k/ [[k]] --- Atividade de Som (não de letra nem de sílaba)]
\noindent\textbf{Nome:} \underline{\hspace{3.2cm}} \quad \textbf{Data:} \underline{\hspace{1.6cm}}

\noindent\textbf{\faMusic~O som /k/ {[k]}:} casa, copo, cubo, queijo, quero, aqui, kiwi, kart.
\newline\textbf{\faPencilRuler~Quem faz esse som:} C, Q, K. \hfill \textbf{\faCheckCircle~Voz:} sem voz.

\vspace{0.4em}
\textbf{A) CAÇA AO SOM} --- Leia cada par com a mão na garganta: marca VIBRA para quem vibra e NÃO VIBRA para quem não vibra.

\begin{center}
\begin{tabular}{@{}ll@{}}
\fbox{\textbf{SIM}} \textit{casa} \quad \fbox{\textbf{NÃO}} \textit{gato} \\
\fbox{\textbf{SIM}} \textit{copo} \quad \fbox{\textbf{NÃO}} \textit{bola} \\
\fbox{\textbf{SIM}} \textit{queijo} \quad \fbox{\textbf{NÃO}} \textit{dedo} \\
\fbox{\textbf{SIM}} \textit{kiwi} \quad \fbox{\textbf{NÃO}} \textit{mala} \\
\end{tabular}
\end{center}

\textbf{B) SEPARE} --- Escreva as palavras na coluna certa.

\begin{center}
\begin{tabular}{@{}p{4.6cm}|p{4.6cm}@{}}
\hline
\textbf{TEM o som} & \textbf{NÃO tem o som} \\
\hline
\textbf{quibe} & \textbf{sapo} \\
\textbf{saco} & \textbf{dedo} \\
\textbf{kart} & \textbf{mesa} \\
\hline
\end{tabular}
\end{center}

\textbf{C) COMPLETE} --- Complete com a grafia do som e confira entre parênteses.

\begin{center}
\large \underline{\hspace{0.7cm}}asa (casa) \quad \underline{\hspace{0.7cm}}opo (copo) \quad \underline{\hspace{0.7cm}}ueijo (queijo) \quad \underline{\hspace{0.7cm}}iwi (kiwi)
\end{center}

\textbf{D) OUÇA E COMPARE} --- K sem voz e G com voz.

\begin{center}
\begin{tabular}{@{}p{4.2cm}p{6.2cm}p{1.6cm}@{}}
\hline
\textbf{Palavra} & \textbf{Compare} & \textbf{Marque} \\
\hline
\textbf{cola} & sem voz [k] & \textbf{NÃO} \\
\textbf{gola} & com voz [g] & \textbf{SIM} \\
\textbf{cata} & sem voz & \textbf{NÃO} \\
\hline
\end{tabular}
\end{center}

\caixapausa{Se precisar, pare aqui. Respire fundo e continue quando quiser.}
\caixapista{Som seco que sai do fundo da boca. Mão na garganta: não vibra (som surdo).}
\end{exercicio}

\plancheta{Atividade de Som 22/38 --- /g/ --- GATO, GOLA e GEMA}
\begin{exercicio}[Som /g/ [[g]] --- Atividade de Som (não de letra nem de sílaba)]
\noindent\textbf{Nome:} \underline{\hspace{3.2cm}} \quad \textbf{Data:} \underline{\hspace{1.6cm}}

\noindent\textbf{\faMusic~O som /g/ {[g]}:} gato, gola, gude; gema, girafa, gelo, página.
\newline\textbf{\faPencilRuler~Quem faz esse som:} G. \hfill \textbf{\faCheckCircle~Voz:} com voz.

\vspace{0.4em}
\textbf{A) CAÇA AO SOM} --- Leia cada par com a mão na garganta: marca VIBRA para quem vibra e NÃO VIBRA para quem não vibra.

\begin{center}
\begin{tabular}{@{}ll@{}}
\fbox{\textbf{SIM}} \textit{gato} \quad \fbox{\textbf{NÃO}} \textit{casa} \\
\fbox{\textbf{SIM}} \textit{gola} \quad \fbox{\textbf{NÃO}} \textit{pato} \\
\fbox{\textbf{SIM}} \textit{gude} \quad \fbox{\textbf{NÃO}} \textit{bola} \\
\fbox{\textbf{SIM}} \textit{gema} \quad \fbox{\textbf{NÃO}} \textit{dedo} \\
\end{tabular}
\end{center}

\textbf{B) SEPARE} --- Escreva as palavras na coluna certa.

\begin{center}
\begin{tabular}{@{}p{4.6cm}|p{4.6cm}@{}}
\hline
\textbf{TEM o som} & \textbf{NÃO tem o som} \\
\hline
\textbf{girafa} & \textbf{sapo} \\
\textbf{gelo} & \textbf{pipa} \\
\textbf{página} & \textbf{teto} \\
\hline
\end{tabular}
\end{center}

\textbf{C) COMPLETE} --- Complete com a grafia do som e confira entre parênteses.

\begin{center}
\large \underline{\hspace{0.7cm}}ato (gato) \quad \underline{\hspace{0.7cm}}ola (gola) \quad \underline{\hspace{0.7cm}}ude (gude) \quad \underline{\hspace{0.7cm}}ema (gema)
\end{center}

\textbf{D) OUÇA E COMPARE} --- G com voz e K sem voz.

\begin{center}
\begin{tabular}{@{}p{4.2cm}p{6.2cm}p{1.6cm}@{}}
\hline
\textbf{Palavra} & \textbf{Compare} & \textbf{Marque} \\
\hline
\textbf{gata} & com voz [g] & \textbf{SIM} \\
\textbf{cata} & sem voz [k] & \textbf{NÃO} \\
\textbf{gola} & com voz & \textbf{SIM} \\
\hline
\end{tabular}
\end{center}

\caixapausa{Se precisar, pare aqui. Respire fundo e continue quando quiser.}
\caixapista{Fundo da boca, COM voz. (G antes de E e I soa como J.) Mão na garganta: vibra.}
\end{exercicio}

\plancheta{Atividade de Som 23/38 --- /f/ --- FACA, FOCA e FIGO}
\begin{exercicio}[Som /f/ [[f]] --- Atividade de Som (não de letra nem de sílaba)]
\noindent\textbf{Nome:} \underline{\hspace{3.2cm}} \quad \textbf{Data:} \underline{\hspace{1.6cm}}

\noindent\textbf{\faMusic~O som /f/ {[f]}:} faca, foca, figo, fumo, fada, café, sofá.
\newline\textbf{\faPencilRuler~Quem faz esse som:} F. \hfill \textbf{\faCheckCircle~Voz:} sem voz.

\vspace{0.4em}
\textbf{A) CAÇA AO SOM} --- Leia cada par com a mão na garganta: marca VIBRA para quem vibra e NÃO VIBRA para quem não vibra.

\begin{center}
\begin{tabular}{@{}ll@{}}
\fbox{\textbf{SIM}} \textit{faca} \quad \fbox{\textbf{NÃO}} \textit{vaca} \\
\fbox{\textbf{SIM}} \textit{foca} \quad \fbox{\textbf{NÃO}} \textit{bola} \\
\fbox{\textbf{SIM}} \textit{figo} \quad \fbox{\textbf{NÃO}} \textit{dedo} \\
\fbox{\textbf{SIM}} \textit{fumo} \quad \fbox{\textbf{NÃO}} \textit{casa} \\
\end{tabular}
\end{center}

\textbf{B) SEPARE} --- Escreva as palavras na coluna certa.

\begin{center}
\begin{tabular}{@{}p{4.6cm}|p{4.6cm}@{}}
\hline
\textbf{TEM o som} & \textbf{NÃO tem o som} \\
\hline
\textbf{sofá} & \textbf{sapo} \\
\textbf{fada} & \textbf{pipa} \\
\textbf{foca} & \textbf{mala} \\
\hline
\end{tabular}
\end{center}

\textbf{C) COMPLETE} --- Complete com a grafia do som e confira entre parênteses.

\begin{center}
\large \underline{\hspace{0.7cm}}aca (faca) \quad \underline{\hspace{0.7cm}}oca (foca) \quad \underline{\hspace{0.7cm}}igo (figo) \quad \underline{\hspace{0.7cm}}umo (fumo)
\end{center}

\textbf{D) OUÇA E COMPARE} --- F sem voz e V com voz.

\begin{center}
\begin{tabular}{@{}p{4.2cm}p{6.2cm}p{1.6cm}@{}}
\hline
\textbf{Palavra} & \textbf{Compare} & \textbf{Marque} \\
\hline
\textbf{faca} & sem voz [f] & \textbf{NÃO} \\
\textbf{vaca} & com voz [v] & \textbf{SIM} \\
\textbf{foto} & sem voz & \textbf{NÃO} \\
\hline
\end{tabular}
\end{center}

\caixapausa{Se precisar, pare aqui. Respire fundo e continue quando quiser.}
\caixapista{Dentes de cima no lábio de baixo; ar forte. Mão na frente da boca: vento frio.}
\end{exercicio}

\plancheta{Atividade de Som 24/38 --- /v/ --- VACA, VELA e VIDA}
\begin{exercicio}[Som /v/ [[v]] --- Atividade de Som (não de letra nem de sílaba)]
\noindent\textbf{Nome:} \underline{\hspace{3.2cm}} \quad \textbf{Data:} \underline{\hspace{1.6cm}}

\noindent\textbf{\faMusic~O som /v/ {[v]}:} vaca, vela, vida, vovô, vulcão, ovo, ave.
\newline\textbf{\faPencilRuler~Quem faz esse som:} V. \hfill \textbf{\faCheckCircle~Voz:} com voz.

\vspace{0.4em}
\textbf{A) CAÇA AO SOM} --- Leia cada par com a mão na garganta: marca VIBRA para quem vibra e NÃO VIBRA para quem não vibra.

\begin{center}
\begin{tabular}{@{}ll@{}}
\fbox{\textbf{SIM}} \textit{vaca} \quad \fbox{\textbf{NÃO}} \textit{faca} \\
\fbox{\textbf{SIM}} \textit{vela} \quad \fbox{\textbf{NÃO}} \textit{bola} \\
\fbox{\textbf{SIM}} \textit{vida} \quad \fbox{\textbf{NÃO}} \textit{dedo} \\
\fbox{\textbf{SIM}} \textit{vovô} \quad \fbox{\textbf{NÃO}} \textit{casa} \\
\end{tabular}
\end{center}

\textbf{B) SEPARE} --- Escreva as palavras na coluna certa.

\begin{center}
\begin{tabular}{@{}p{4.6cm}|p{4.6cm}@{}}
\hline
\textbf{TEM o som} & \textbf{NÃO tem o som} \\
\hline
\textbf{ovo} & \textbf{sapo} \\
\textbf{ave} & \textbf{pipa} \\
\textbf{vulcão} & \textbf{mala} \\
\hline
\end{tabular}
\end{center}

\textbf{C) COMPLETE} --- Complete com a grafia do som e confira entre parênteses.

\begin{center}
\large \underline{\hspace{0.7cm}}aca (vaca) \quad \underline{\hspace{0.7cm}}ela (vela) \quad \underline{\hspace{0.7cm}}ida (vida) \quad \underline{\hspace{0.7cm}}ovô (vovô)
\end{center}

\textbf{D) OUÇA E COMPARE} --- V com voz e F sem voz.

\begin{center}
\begin{tabular}{@{}p{4.2cm}p{6.2cm}p{1.6cm}@{}}
\hline
\textbf{Palavra} & \textbf{Compare} & \textbf{Marque} \\
\hline
\textbf{vaca} & com voz [v] & \textbf{SIM} \\
\textbf{faca} & sem voz [f] & \textbf{NÃO} \\
\textbf{voto} & com voz & \textbf{SIM} \\
\hline
\end{tabular}
\end{center}

\caixapausa{Se precisar, pare aqui. Respire fundo e continue quando quiser.}
\caixapista{Dentes de cima no lábio de baixo, COM voz. Mão na garganta: vibra.}
\end{exercicio}

\plancheta{Atividade de Som 25/38 --- /s/ --- SAPO, CEDO e PAZ}
\begin{exercicio}[Som /s/ [[s]] --- Atividade de Som (não de letra nem de sílaba)]
\noindent\textbf{Nome:} \underline{\hspace{3.2cm}} \quad \textbf{Data:} \underline{\hspace{1.6cm}}

\noindent\textbf{\faMusic~O som /s/ {[s]}:} sapo, sopa, sala, sino; cedo, céu, doce; caça, poço; paz, luz; próximo.
\newline\textbf{\faPencilRuler~Quem faz esse som:} S, C, Ç, X, Z. \hfill \textbf{\faCheckCircle~Voz:} sem voz.

\vspace{0.4em}
\textbf{A) CAÇA AO SOM} --- Leia cada par com a mão na garganta: marca VIBRA para quem vibra e NÃO VIBRA para quem não vibra.

\begin{center}
\begin{tabular}{@{}ll@{}}
\fbox{\textbf{SIM}} \textit{sapo} \quad \fbox{\textbf{NÃO}} \textit{casa} \\
\fbox{\textbf{SIM}} \textit{sopa} \quad \fbox{\textbf{NÃO}} \textit{bola} \\
\fbox{\textbf{SIM}} \textit{sino} \quad \fbox{\textbf{NÃO}} \textit{dedo} \\
\fbox{\textbf{SIM}} \textit{doce} \quad \fbox{\textbf{NÃO}} \textit{gato} \\
\end{tabular}
\end{center}

\textbf{B) SEPARE} --- Escreva as palavras na coluna certa.

\begin{center}
\begin{tabular}{@{}p{4.6cm}|p{4.6cm}@{}}
\hline
\textbf{TEM o som} & \textbf{NÃO tem o som} \\
\hline
\textbf{caça} & \textbf{mala} \\
\textbf{poço} & \textbf{pipa} \\
\textbf{luz} & \textbf{teto} \\
\hline
\end{tabular}
\end{center}

\textbf{C) COMPLETE} --- Complete com a grafia do som e confira entre parênteses.

\begin{center}
\large \underline{\hspace{0.8cm}}apo (sapo) \quad \underline{\hspace{0.8cm}}edo (c) (cedo) \quad \underline{\hspace{0.8cm}}aça (ç) (caça) \quad pa\underline{\hspace{0.8cm}} (z) (paz)
\end{center}

\textbf{D) OUÇA E COMPARE} --- S sem voz e Z com voz.

\begin{center}
\begin{tabular}{@{}p{4.2cm}p{6.2cm}p{1.6cm}@{}}
\hline
\textbf{Palavra} & \textbf{Compare} & \textbf{Marque} \\
\hline
\textbf{selo} & sem voz [s] & \textbf{NÃO} \\
\textbf{zelo} & com voz [z] & \textbf{SIM} \\
\textbf{caça} & sem voz & \textbf{NÃO} \\
\hline
\end{tabular}
\end{center}

\caixapausa{Se precisar, pare aqui. Respire fundo e continue quando quiser.}
\caixapista{Dentes quase juntos; ar fino e contínuo. Mão na frente da boca: ar frio.}
\end{exercicio}

\plancheta{Atividade de Som 26/38 --- /z/ --- ZEBRA, CASA e MESA}
\begin{exercicio}[Som /z/ [[z]] --- Atividade de Som (não de letra nem de sílaba)]
\noindent\textbf{Nome:} \underline{\hspace{3.2cm}} \quad \textbf{Data:} \underline{\hspace{1.6cm}}

\noindent\textbf{\faMusic~O som /z/ {[z]}:} zero, zebra, zelo; casa, mesa, asa; exame, exemplo.
\newline\textbf{\faPencilRuler~Quem faz esse som:} Z, S, X. \hfill \textbf{\faCheckCircle~Voz:} com voz.

\vspace{0.4em}
\textbf{A) CAÇA AO SOM} --- Leia cada par com a mão na garganta: marca VIBRA para quem vibra e NÃO VIBRA para quem não vibra.

\begin{center}
\begin{tabular}{@{}ll@{}}
\fbox{\textbf{SIM}} \textit{zero} \quad \fbox{\textbf{NÃO}} \textit{sapo} \\
\fbox{\textbf{SIM}} \textit{zebra} \quad \fbox{\textbf{NÃO}} \textit{bola} \\
\fbox{\textbf{SIM}} \textit{casa} \quad \fbox{\textbf{NÃO}} \textit{dedo} \\
\fbox{\textbf{SIM}} \textit{mesa} \quad \fbox{\textbf{NÃO}} \textit{gato} \\
\end{tabular}
\end{center}

\textbf{B) SEPARE} --- Escreva as palavras na coluna certa.

\begin{center}
\begin{tabular}{@{}p{4.6cm}|p{4.6cm}@{}}
\hline
\textbf{TEM o som} & \textbf{NÃO tem o som} \\
\hline
\textbf{exame} & \textbf{mala} \\
\textbf{zíper} & \textbf{pipa} \\
\textbf{asa} & \textbf{teto} \\
\hline
\end{tabular}
\end{center}

\textbf{C) COMPLETE} --- Complete com a grafia do som e confira entre parênteses.

\begin{center}
\large \underline{\hspace{0.8cm}}ero (z) (zero) \quad \underline{\hspace{0.8cm}}ebra (z) (zebra) \quad ca\underline{\hspace{0.7cm}}a (s) (casa) \quad me\underline{\hspace{0.7cm}}a (s) (mesa)
\end{center}

\textbf{D) OUÇA E COMPARE} --- Z com voz e S sem voz.

\begin{center}
\begin{tabular}{@{}p{4.2cm}p{6.2cm}p{1.6cm}@{}}
\hline
\textbf{Palavra} & \textbf{Compare} & \textbf{Marque} \\
\hline
\textbf{zelo} & com voz [z] & \textbf{SIM} \\
\textbf{selo} & sem voz [s] & \textbf{NÃO} \\
\textbf{casa} & com voz & \textbf{SIM} \\
\hline
\end{tabular}
\end{center}

\caixapausa{Se precisar, pare aqui. Respire fundo e continue quando quiser.}
\caixapista{Dentes quase juntos, COM voz (zumbido de abelha). Mão na garganta: vibra.}
\end{exercicio}

\plancheta{Atividade de Som 27/38 --- /ch/ (x) --- XÍCARA, CHUVA e PEIXE}
\begin{exercicio}[Som /ch/ (x) [[S]] --- Atividade de Som (não de letra nem de sílaba)]
\noindent\textbf{Nome:} \underline{\hspace{3.2cm}} \quad \textbf{Data:} \underline{\hspace{1.6cm}}

\noindent\textbf{\faMusic~O som /ch/ (x) {[S]}:} xícara, xale, caixa, peixe; chuva, chave, chão, bicho.
\newline\textbf{\faPencilRuler~Quem faz esse som:} X, CH. \hfill \textbf{\faCheckCircle~Voz:} sem voz.

\vspace{0.4em}
\textbf{A) CAÇA AO SOM} --- Leia cada par com a mão na garganta: marca VIBRA para quem vibra e NÃO VIBRA para quem não vibra.

\begin{center}
\begin{tabular}{@{}ll@{}}
\fbox{\textbf{SIM}} \textit{xícara} \quad \fbox{\textbf{NÃO}} \textit{sapo} \\
\fbox{\textbf{SIM}} \textit{chave} \quad \fbox{\textbf{NÃO}} \textit{bola} \\
\fbox{\textbf{SIM}} \textit{caixa} \quad \fbox{\textbf{NÃO}} \textit{dedo} \\
\fbox{\textbf{SIM}} \textit{peixe} \quad \fbox{\textbf{NÃO}} \textit{gato} \\
\end{tabular}
\end{center}

\textbf{B) SEPARE} --- Escreva as palavras na coluna certa.

\begin{center}
\begin{tabular}{@{}p{4.6cm}|p{4.6cm}@{}}
\hline
\textbf{TEM o som} & \textbf{NÃO tem o som} \\
\hline
\textbf{bicho} & \textbf{mala} \\
\textbf{chão} & \textbf{pipa} \\
\textbf{tacho} & \textbf{teto} \\
\hline
\end{tabular}
\end{center}

\textbf{C) COMPLETE} --- Complete com a grafia do som e confira entre parênteses.

\begin{center}
\large \underline{\hspace{0.8cm}}ícara (x) (xícara) \quad \underline{\hspace{0.8cm}}uva (ch) (chuva) \quad \underline{\hspace{0.8cm}}ave (ch) (chave) \quad cai\underline{\hspace{0.7cm}}a (x) (caixa)
\end{center}

\textbf{D) OUÇA E COMPARE} --- X/CH e J.

\begin{center}
\begin{tabular}{@{}p{4.2cm}p{6.2cm}p{1.6cm}@{}}
\hline
\textbf{Palavra} & \textbf{Compare} & \textbf{Marque} \\
\hline
\textbf{chato} & som de ch [ch] & \textbf{NÃO} \\
\textbf{jato} & som de j [j] & \textbf{SIM} \\
\textbf{chuva} & ch [ch] & \textbf{NÃO} \\
\hline
\end{tabular}
\end{center}

\caixapausa{Se precisar, pare aqui. Respire fundo e continue quando quiser.}
\caixapista{Língua no céu da boca; ar forte, sem voz. Mão na frente da boca: ar forte.}
\end{exercicio}

\plancheta{Atividade de Som 28/38 --- /j/ --- JACA, JOGO e GELO}
\begin{exercicio}[Som /j/ [[Z]] --- Atividade de Som (não de letra nem de sílaba)]
\noindent\textbf{Nome:} \underline{\hspace{3.2cm}} \quad \textbf{Data:} \underline{\hspace{1.6cm}}

\noindent\textbf{\faMusic~O som /j/ {[Z]}:} jaca, jogo, jipe, juba, jejum; gelo, girafa, gema, página.
\newline\textbf{\faPencilRuler~Quem faz esse som:} J, G (e, i). \hfill \textbf{\faCheckCircle~Voz:} com voz.

\vspace{0.4em}
\textbf{A) CAÇA AO SOM} --- Leia cada par com a mão na garganta: marca VIBRA para quem vibra e NÃO VIBRA para quem não vibra.

\begin{center}
\begin{tabular}{@{}ll@{}}
\fbox{\textbf{SIM}} \textit{jaca} \quad \fbox{\textbf{NÃO}} \textit{chave} \\
\fbox{\textbf{SIM}} \textit{jogo} \quad \fbox{\textbf{NÃO}} \textit{bola} \\
\fbox{\textbf{SIM}} \textit{gelo} \quad \fbox{\textbf{NÃO}} \textit{dedo} \\
\fbox{\textbf{SIM}} \textit{juba} \quad \fbox{\textbf{NÃO}} \textit{casa} \\
\end{tabular}
\end{center}

\textbf{B) SEPARE} --- Escreva as palavras na coluna certa.

\begin{center}
\begin{tabular}{@{}p{4.6cm}|p{4.6cm}@{}}
\hline
\textbf{TEM o som} & \textbf{NÃO tem o som} \\
\hline
\textbf{página} & \textbf{mala} \\
\textbf{jejum} & \textbf{pipa} \\
\textbf{tigela} & \textbf{teto} \\
\hline
\end{tabular}
\end{center}

\textbf{C) COMPLETE} --- Complete com a grafia do som e confira entre parênteses.

\begin{center}
\large \underline{\hspace{0.7cm}}aca (j) (jaca) \quad \underline{\hspace{0.7cm}}ogo (j) (jogo) \quad \underline{\hspace{0.7cm}}ube (j) (jube (leia devagar)) \quad \underline{\hspace{0.7cm}}ema (g) (gema)
\end{center}

\textbf{D) OUÇA E COMPARE} --- J e X/CH.

\begin{center}
\begin{tabular}{@{}p{4.2cm}p{6.2cm}p{1.6cm}@{}}
\hline
\textbf{Palavra} & \textbf{Compare} & \textbf{Marque} \\
\hline
\textbf{jato} & com voz [j] & \textbf{SIM} \\
\textbf{chato} & sem voz [ch] & \textbf{NÃO} \\
\textbf{juba} & com voz & \textbf{SIM} \\
\hline
\end{tabular}
\end{center}

\caixapausa{Se precisar, pare aqui. Respire fundo e continue quando quiser.}
\caixapista{Língua no céu da boca, COM voz (som de J). Mão na garganta: vibra.}
\end{exercicio}

\plancheta{Atividade de Som 29/38 --- /m/ --- MALA, MÃO e MAR}
\begin{exercicio}[Som /m/ [[m]] --- Atividade de Som (não de letra nem de sílaba)]
\noindent\textbf{Nome:} \underline{\hspace{3.2cm}} \quad \textbf{Data:} \underline{\hspace{1.6cm}}

\noindent\textbf{\faMusic~O som /m/ {[m]}:} mala, mão, mim, mar, mel; (fim) bom, som, campo.
\newline\textbf{\faPencilRuler~Quem faz esse som:} M. \hfill \textbf{\faCheckCircle~Voz:} com voz, nasal.

\vspace{0.4em}
\textbf{A) CAÇA AO SOM} --- Leia cada par com o dedo no nariz: marca NARIZ para quem faz o nariz vibrar e BOCA para quem não faz.

\begin{center}
\begin{tabular}{@{}ll@{}}
\fbox{\textbf{SIM}} \textit{mala} \quad \fbox{\textbf{NÃO}} \textit{nada} \\
\fbox{\textbf{SIM}} \textit{mão} \quad \fbox{\textbf{NÃO}} \textit{bola} \\
\fbox{\textbf{SIM}} \textit{mim} \quad \fbox{\textbf{NÃO}} \textit{dedo} \\
\fbox{\textbf{SIM}} \textit{mar} \quad \fbox{\textbf{NÃO}} \textit{casa} \\
\end{tabular}
\end{center}

\textbf{B) SEPARE} --- Escreva as palavras na coluna certa.

\begin{center}
\begin{tabular}{@{}p{4.6cm}|p{4.6cm}@{}}
\hline
\textbf{TEM o som} & \textbf{NÃO tem o som} \\
\hline
\textbf{mel} & \textbf{sapo} \\
\textbf{mula} & \textbf{pipa} \\
\textbf{campo} & \textbf{teto} \\
\hline
\end{tabular}
\end{center}

\textbf{C) COMPLETE} --- Complete com a grafia do som e confira entre parênteses.

\begin{center}
\large \underline{\hspace{0.7cm}}ala (mala) \quad \underline{\hspace{0.7cm}}ão (mão) \quad \underline{\hspace{0.7cm}}im (mim) \quad \underline{\hspace{0.7cm}}ar (mar)
\end{center}

\textbf{D) OUÇA E COMPARE} --- M e N.

\begin{center}
\begin{tabular}{@{}p{4.2cm}p{6.2cm}p{1.6cm}@{}}
\hline
\textbf{Palavra} & \textbf{Compare} & \textbf{Marque} \\
\hline
\textbf{mão} & começa com M [m] & \textbf{NÃO} \\
\textbf{não} & começa com N [n] & \textbf{SIM} \\
\textbf{mula} & M & \textbf{NÃO} \\
\hline
\end{tabular}
\end{center}

\caixapausa{Se precisar, pare aqui. Respire fundo e continue quando quiser.}
\caixapista{Boca fechada; ar sai pelo nariz. Dedo no nariz: sente a vibração nasal.}
\end{exercicio}

\plancheta{Atividade de Som 30/38 --- /n/ --- NADA, NENA e NUVEM}
\begin{exercicio}[Som /n/ [[n]] --- Atividade de Som (não de letra nem de sílaba)]
\noindent\textbf{Nome:} \underline{\hspace{3.2cm}} \quad \textbf{Data:} \underline{\hspace{1.6cm}}

\noindent\textbf{\faMusic~O som /n/ {[n]}:} nada, nena, nó, nuvem, cana, sino.
\newline\textbf{\faPencilRuler~Quem faz esse som:} N. \hfill \textbf{\faCheckCircle~Voz:} com voz, nasal.

\vspace{0.4em}
\textbf{A) CAÇA AO SOM} --- Leia cada par com o dedo no nariz: marca NARIZ para quem faz o nariz vibrar e BOCA para quem não faz.

\begin{center}
\begin{tabular}{@{}ll@{}}
\fbox{\textbf{SIM}} \textit{nada} \quad \fbox{\textbf{NÃO}} \textit{mala} \\
\fbox{\textbf{SIM}} \textit{nena} \quad \fbox{\textbf{NÃO}} \textit{bola} \\
\fbox{\textbf{SIM}} \textit{cana} \quad \fbox{\textbf{NÃO}} \textit{dedo} \\
\fbox{\textbf{SIM}} \textit{nó} \quad \fbox{\textbf{NÃO}} \textit{casa} \\
\end{tabular}
\end{center}

\textbf{B) SEPARE} --- Escreva as palavras na coluna certa.

\begin{center}
\begin{tabular}{@{}p{4.6cm}|p{4.6cm}@{}}
\hline
\textbf{TEM o som} & \textbf{NÃO tem o som} \\
\hline
\textbf{nuvem} & \textbf{sapo} \\
\textbf{ninho} & \textbf{pipa} \\
\textbf{sino} & \textbf{teto} \\
\hline
\end{tabular}
\end{center}

\textbf{C) COMPLETE} --- Complete com a grafia do som e confira entre parênteses.

\begin{center}
\large \underline{\hspace{0.7cm}}ada (nada) \quad \underline{\hspace{0.7cm}}ena (nena) \quad \underline{\hspace{0.7cm}}uvem (nuvem) \quad ca\underline{\hspace{0.7cm}}a (cana)
\end{center}

\textbf{D) OUÇA E COMPARE} --- N e M.

\begin{center}
\begin{tabular}{@{}p{4.2cm}p{6.2cm}p{1.6cm}@{}}
\hline
\textbf{Palavra} & \textbf{Compare} & \textbf{Marque} \\
\hline
\textbf{não} & N [n] & \textbf{SIM} \\
\textbf{mão} & M [m] & \textbf{NÃO} \\
\textbf{nula} & N & \textbf{SIM} \\
\hline
\end{tabular}
\end{center}

\caixapausa{Se precisar, pare aqui. Respire fundo e continue quando quiser.}
\caixapista{Língua no céu da boca; ar pelo nariz. Dedo no nariz: vibra.}
\end{exercicio}

\plancheta{Atividade de Som 31/38 --- /nh/ --- NINHO, UNHA e SONHO}
\begin{exercicio}[Som /nh/ [[\textltailn]] --- Atividade de Som (não de letra nem de sílaba)]
\noindent\textbf{Nome:} \underline{\hspace{3.2cm}} \quad \textbf{Data:} \underline{\hspace{1.6cm}}

\noindent\textbf{\faMusic~O som /nh/ {[\textltailn]}:} ninho, unha, sonho, banho, linha, vinho, aranha.
\newline\textbf{\faPencilRuler~Quem faz esse som:} NH. \hfill \textbf{\faCheckCircle~Voz:} com voz, nasal.

\vspace{0.4em}
\textbf{A) CAÇA AO SOM} --- Leia cada par com o dedo no nariz: marca NARIZ para quem faz o nariz vibrar e BOCA para quem não faz.

\begin{center}
\begin{tabular}{@{}ll@{}}
\fbox{\textbf{SIM}} \textit{ninho} \quad \fbox{\textbf{NÃO}} \textit{nada} \\
\fbox{\textbf{SIM}} \textit{unha} \quad \fbox{\textbf{NÃO}} \textit{bola} \\
\fbox{\textbf{SIM}} \textit{sonho} \quad \fbox{\textbf{NÃO}} \textit{dedo} \\
\fbox{\textbf{SIM}} \textit{banho} \quad \fbox{\textbf{NÃO}} \textit{casa} \\
\end{tabular}
\end{center}

\textbf{B) SEPARE} --- Escreva as palavras na coluna certa.

\begin{center}
\begin{tabular}{@{}p{4.6cm}|p{4.6cm}@{}}
\hline
\textbf{TEM o som} & \textbf{NÃO tem o som} \\
\hline
\textbf{linha} & \textbf{sapo} \\
\textbf{vinho} & \textbf{pipa} \\
\textbf{aranha} & \textbf{teto} \\
\hline
\end{tabular}
\end{center}

\textbf{C) COMPLETE} --- Complete com a grafia do som e confira entre parênteses.

\begin{center}
\large ni\underline{\hspace{0.7cm}}ho (nh) (ninho) \quad u\underline{\hspace{0.7cm}}a (nh) (unha) \quad so\underline{\hspace{0.7cm}}o (nh) (sonho) \quad ba\underline{\hspace{0.7cm}}o (nh) (banho)
\end{center}

\textbf{D) OUÇA E COMPARE} --- NH e N.

\begin{center}
\begin{tabular}{@{}p{4.2cm}p{6.2cm}p{1.6cm}@{}}
\hline
\textbf{Palavra} & \textbf{Compare} & \textbf{Marque} \\
\hline
\textbf{ninho} & NH no meio [nh] & \textbf{NÃO} \\
\textbf{nina (leia devagar)} & N [n] & \textbf{SIM} \\
\textbf{unha} & NH [nh] & \textbf{NÃO} \\
\hline
\end{tabular}
\end{center}

\caixapausa{Se precisar, pare aqui. Respire fundo e continue quando quiser.}
\caixapista{Língua no céu da boca; ar pelo nariz (NH molhado). Dedo no nariz: vibra.}
\end{exercicio}

\plancheta{Atividade de Som 32/38 --- /lh/ --- MILHO, ALHO e FILHO}
\begin{exercicio}[Som /lh/ [[\textturny]] --- Atividade de Som (não de letra nem de sílaba)]
\noindent\textbf{Nome:} \underline{\hspace{3.2cm}} \quad \textbf{Data:} \underline{\hspace{1.6cm}}

\noindent\textbf{\faMusic~O som /lh/ {[\textturny]}:} milho, alho, filho, olho, ilha, coelho.
\newline\textbf{\faPencilRuler~Quem faz esse som:} LH. \hfill \textbf{\faCheckCircle~Voz:} com voz.

\vspace{0.4em}
\textbf{A) CAÇA AO SOM} --- Leia em voz alta e compare os sons. Marque a coluna certa para cada palavra.

\begin{center}
\begin{tabular}{@{}ll@{}}
\fbox{\textbf{SIM}} \textit{milho} \quad \fbox{\textbf{NÃO}} \textit{mala} \\
\fbox{\textbf{SIM}} \textit{alho} \quad \fbox{\textbf{NÃO}} \textit{bola} \\
\fbox{\textbf{SIM}} \textit{filho} \quad \fbox{\textbf{NÃO}} \textit{dedo} \\
\fbox{\textbf{SIM}} \textit{olho} \quad \fbox{\textbf{NÃO}} \textit{casa} \\
\end{tabular}
\end{center}

\textbf{B) SEPARE} --- Escreva as palavras na coluna certa.

\begin{center}
\begin{tabular}{@{}p{4.6cm}|p{4.6cm}@{}}
\hline
\textbf{TEM o som} & \textbf{NÃO tem o som} \\
\hline
\textbf{ilha} & \textbf{sapo} \\
\textbf{coelho} & \textbf{pipa} \\
\textbf{baralho} & \textbf{teto} \\
\hline
\end{tabular}
\end{center}

\textbf{C) COMPLETE} --- Complete com a grafia do som e confira entre parênteses.

\begin{center}
\large mi\underline{\hspace{0.7cm}}o (lh) (milho) \quad a\underline{\hspace{0.7cm}}o (lh) (alho) \quad fi\underline{\hspace{0.7cm}}o (lh) (filho) \quad o\underline{\hspace{0.7cm}}o (lh) (olho)
\end{center}

\textbf{D) OUÇA E COMPARE} --- LH e L.

\begin{center}
\begin{tabular}{@{}p{4.2cm}p{6.2cm}p{1.6cm}@{}}
\hline
\textbf{Palavra} & \textbf{Compare} & \textbf{Marque} \\
\hline
\textbf{milho} & LH [lh] & \textbf{NÃO} \\
\textbf{filó (leia devagar)} & L [l] & \textbf{SIM} \\
\textbf{olho} & LH & \textbf{NÃO} \\
\hline
\end{tabular}
\end{center}

\caixapausa{Se precisar, pare aqui. Respire fundo e continue quando quiser.}
\caixapista{Língua no céu da boca; som molhado (LH). Mão na lateral do pescoço: vibra.}
\end{exercicio}

\plancheta{Atividade de Som 33/38 --- /l/ --- LATA, LUA e SALA}
\begin{exercicio}[Som /l/ [[l]] --- Atividade de Som (não de letra nem de sílaba)]
\noindent\textbf{Nome:} \underline{\hspace{3.2cm}} \quad \textbf{Data:} \underline{\hspace{1.6cm}}

\noindent\textbf{\faMusic~O som /l/ {[l]}:} lata, lua, leão, sala, bala, lápis.
\newline\textbf{\faPencilRuler~Quem faz esse som:} L. \hfill \textbf{\faCheckCircle~Voz:} com voz.

\vspace{0.4em}
\textbf{A) CAÇA AO SOM} --- Leia em voz alta e compare os sons. Marque a coluna certa para cada palavra.

\begin{center}
\begin{tabular}{@{}ll@{}}
\fbox{\textbf{SIM}} \textit{lata} \quad \fbox{\textbf{NÃO}} \textit{rua} \\
\fbox{\textbf{SIM}} \textit{lua} \quad \fbox{\textbf{NÃO}} \textit{bola} \\
\fbox{\textbf{SIM}} \textit{leão} \quad \fbox{\textbf{NÃO}} \textit{dedo} \\
\fbox{\textbf{SIM}} \textit{sala} \quad \fbox{\textbf{NÃO}} \textit{casa} \\
\end{tabular}
\end{center}

\textbf{B) SEPARE} --- Escreva as palavras na coluna certa.

\begin{center}
\begin{tabular}{@{}p{4.6cm}|p{4.6cm}@{}}
\hline
\textbf{TEM o som} & \textbf{NÃO tem o som} \\
\hline
\textbf{balão} & \textbf{sapo} \\
\textbf{lápis} & \textbf{pipa} \\
\textbf{lindo} & \textbf{teto} \\
\hline
\end{tabular}
\end{center}

\textbf{C) COMPLETE} --- Complete com a grafia do som e confira entre parênteses.

\begin{center}
\large \underline{\hspace{0.7cm}}ata (lata) \quad \underline{\hspace{0.7cm}}ua (lua) \quad sa\underline{\hspace{0.7cm}}a (sala) \quad ba\underline{\hspace{0.7cm}}a (bala)
\end{center}

\textbf{D) OUÇA E COMPARE} --- L e R fraco.

\begin{center}
\begin{tabular}{@{}p{4.2cm}p{6.2cm}p{1.6cm}@{}}
\hline
\textbf{Palavra} & \textbf{Compare} & \textbf{Marque} \\
\hline
\textbf{lata} & L [l] & \textbf{NÃO} \\
\textbf{rata} & R [r fraco] & \textbf{SIM} \\
\textbf{lua} & L & \textbf{NÃO} \\
\hline
\end{tabular}
\end{center}

\caixapausa{Se precisar, pare aqui. Respire fundo e continue quando quiser.}
\caixapista{Ponta da língua no céu da boca; ar sai pelos lados. Mão na lateral do pescoço: vibra (L no fim vira [w] em SOL, MEL).}
\end{exercicio}

\plancheta{Atividade de Som 34/38 --- /r/ fraco --- CARA, PERA e ARO}
\begin{exercicio}[Som /r/ fraco [[r]] --- Atividade de Som (não de letra nem de sílaba)]
\noindent\textbf{Nome:} \underline{\hspace{3.2cm}} \quad \textbf{Data:} \underline{\hspace{1.6cm}}

\noindent\textbf{\faMusic~O som /r/ fraco {[r]}:} cara, pera, aro, barato, coração.
\newline\textbf{\faPencilRuler~Quem faz esse som:} R fraco. \hfill \textbf{\faCheckCircle~Voz:} com voz.

\vspace{0.4em}
\textbf{A) CAÇA AO SOM} --- Leia em voz alta e compare os sons. Marque a coluna certa para cada palavra.

\begin{center}
\begin{tabular}{@{}ll@{}}
\fbox{\textbf{SIM}} \textit{cara} \quad \fbox{\textbf{NÃO}} \textit{rua} \\
\fbox{\textbf{SIM}} \textit{pera} \quad \fbox{\textbf{NÃO}} \textit{bola} \\
\fbox{\textbf{SIM}} \textit{aro} \quad \fbox{\textbf{NÃO}} \textit{dedo} \\
\fbox{\textbf{SIM}} \textit{barato} \quad \fbox{\textbf{NÃO}} \textit{casa} \\
\end{tabular}
\end{center}

\textbf{B) SEPARE} --- Escreva as palavras na coluna certa.

\begin{center}
\begin{tabular}{@{}p{4.6cm}|p{4.6cm}@{}}
\hline
\textbf{TEM o som} & \textbf{NÃO tem o som} \\
\hline
\textbf{coro} & \textbf{sapo} \\
\textbf{sereia} & \textbf{pipa} \\
\textbf{carinho} & \textbf{teto} \\
\hline
\end{tabular}
\end{center}

\textbf{C) COMPLETE} --- Complete com a grafia do som e confira entre parênteses.

\begin{center}
\large ca\underline{\hspace{0.7cm}}a (cara) \quad pe\underline{\hspace{0.7cm}}a (pera) \quad a\underline{\hspace{0.7cm}}o (aro) \quad ba\underline{\hspace{0.8cm}}ato (barato)
\end{center}

\textbf{D) OUÇA E COMPARE} --- R fraco e R forte.

\begin{center}
\begin{tabular}{@{}p{4.2cm}p{6.2cm}p{1.6cm}@{}}
\hline
\textbf{Palavra} & \textbf{Compare} & \textbf{Marque} \\
\hline
\textbf{cara} & uma batidinha [r fraco] & \textbf{NÃO} \\
\textbf{carro} & som forte [h] & \textbf{SIM} \\
\textbf{coro} & fraco & \textbf{NÃO} \\
\hline
\end{tabular}
\end{center}

\caixapausa{Se precisar, pare aqui. Respire fundo e continue quando quiser.}
\caixapista{Língua dá uma batidinha leve (R fraco). Mão na garganta: vibra.}
\end{exercicio}

\plancheta{Atividade de Som 35/38 --- /h/ --- RATO, RUA e CARRO}
\begin{exercicio}[Som /h/ [[h]] --- Atividade de Som (não de letra nem de sílaba)]
\noindent\textbf{Nome:} \underline{\hspace{3.2cm}} \quad \textbf{Data:} \underline{\hspace{1.6cm}}

\noindent\textbf{\faMusic~O som /h/ {[h]}:} rato, rua, rede, rosa, carro; mar, amor, flor.
\newline\textbf{\faPencilRuler~Quem faz esse som:} R forte. \hfill \textbf{\faCheckCircle~Voz:} com voz.

\vspace{0.4em}
\textbf{A) CAÇA AO SOM} --- Leia em voz alta e compare os sons. Marque a coluna certa para cada palavra.

\begin{center}
\begin{tabular}{@{}ll@{}}
\fbox{\textbf{SIM}} \textit{rato} \quad \fbox{\textbf{NÃO}} \textit{cara} \\
\fbox{\textbf{SIM}} \textit{rua} \quad \fbox{\textbf{NÃO}} \textit{bola} \\
\fbox{\textbf{SIM}} \textit{mar} \quad \fbox{\textbf{NÃO}} \textit{dedo} \\
\fbox{\textbf{SIM}} \textit{carro} \quad \fbox{\textbf{NÃO}} \textit{casa} \\
\end{tabular}
\end{center}

\textbf{B) SEPARE} --- Escreva as palavras na coluna certa.

\begin{center}
\begin{tabular}{@{}p{4.6cm}|p{4.6cm}@{}}
\hline
\textbf{TEM o som} & \textbf{NÃO tem o som} \\
\hline
\textbf{rosa} & \textbf{sapo} \\
\textbf{rede} & \textbf{pipa} \\
\textbf{amor} & \textbf{teto} \\
\hline
\end{tabular}
\end{center}

\textbf{C) COMPLETE} --- Complete com a grafia do som e confira entre parênteses.

\begin{center}
\large \underline{\hspace{0.7cm}}ato (rato) \quad \underline{\hspace{0.7cm}}ua (rua) \quad \underline{\hspace{0.7cm}}osa (rosa) \quad \underline{\hspace{0.7cm}}em (rem)
\end{center}

\textbf{D) OUÇA E COMPARE} --- R forte e R fraco.

\begin{center}
\begin{tabular}{@{}p{4.2cm}p{6.2cm}p{1.6cm}@{}}
\hline
\textbf{Palavra} & \textbf{Compare} & \textbf{Marque} \\
\hline
\textbf{carro} & forte [h] & \textbf{SIM} \\
\textbf{cara} & fraco [r fraco] & \textbf{NÃO} \\
\textbf{corro} & forte & \textbf{SIM} \\
\hline
\end{tabular}
\end{center}

\caixapausa{Se precisar, pare aqui. Respire fundo e continue quando quiser.}
\caixapista{Som forte do fundo da garganta (RR ou H, conforme a região). Mão na garganta: vibra forte.}
\end{exercicio}

\plancheta{Atividade de Som 36/38 --- /tch/ --- TIA, LEITE e TCHAU}
\begin{exercicio}[Som /tch/ [[tS]] --- Atividade de Som (não de letra nem de sílaba)]
\noindent\textbf{Nome:} \underline{\hspace{3.2cm}} \quad \textbf{Data:} \underline{\hspace{1.6cm}}

\noindent\textbf{\faMusic~O som /tch/ {[tS]}:} tia, leite, noite, tchau, quente, gente (variação regional).
\newline\textbf{\faPencilRuler~Quem faz esse som:} T + I (regional). \hfill \textbf{\faCheckCircle~Voz:} sem voz (variação regional).

\vspace{0.4em}
\textbf{A) CAÇA AO SOM} --- Em muitas regiões do Brasil a pronúncia varia. Leia devagar e compare: você fala TCHIA ou TIA? DJIA ou DIA?

\begin{center}
\begin{tabular}{@{}ll@{}}
\fbox{\textbf{SIM}} \textit{tchau} \quad \fbox{\textbf{NÃO}} \textit{teto} \\
\fbox{\textbf{SIM}} \textit{leite} \quad \fbox{\textbf{NÃO}} \textit{bola} \\
\fbox{\textbf{SIM}} \textit{tia} \quad \fbox{\textbf{NÃO}} \textit{dedo} \\
\fbox{\textbf{SIM}} \textit{noite} \quad \fbox{\textbf{NÃO}} \textit{casa} \\
\end{tabular}
\end{center}

\textbf{B) SEPARE} --- Escreva as palavras na coluna certa.

\begin{center}
\begin{tabular}{@{}p{4.6cm}|p{4.6cm}@{}}
\hline
\textbf{TEM o som} & \textbf{NÃO tem o som} \\
\hline
\textbf{gente} & \textbf{sapo} \\
\textbf{quente} & \textbf{pipa} \\
\textbf{leite} & \textbf{teto} \\
\hline
\end{tabular}
\end{center}

\textbf{C) COMPLETE} --- Complete com a grafia do som e confira entre parênteses.

\begin{center}
\large \underline{\hspace{0.7cm}}ia (tia) \quad lei\underline{\hspace{0.7cm}}e (leite) \quad noi\underline{\hspace{0.7cm}}e (noite) \quad \underline{\hspace{0.8cm}}chau (tchau)
\end{center}

\textbf{D) OUÇA E COMPARE} --- T seco e T de TIA.

\begin{center}
\begin{tabular}{@{}p{4.2cm}p{6.2cm}p{1.6cm}@{}}
\hline
\textbf{Palavra} & \textbf{Compare} & \textbf{Marque} \\
\hline
\textbf{teto} & T seco [t] & \textbf{NÃO} \\
\textbf{tia} & T + I vira TCH [tch] & \textbf{SIM} \\
\textbf{pote} & seco & \textbf{NÃO} \\
\hline
\end{tabular}
\end{center}

\caixapausa{Se precisar, pare aqui. Respire fundo e continue quando quiser.}
\caixapista{Língua encosta e solta de repente (TCH). Mão na garganta: não vibra.}
\end{exercicio}

\plancheta{Atividade de Som 37/38 --- /dj/ --- DIA, DICA e DENTE}
\begin{exercicio}[Som /dj/ [[dZ]] --- Atividade de Som (não de letra nem de sílaba)]
\noindent\textbf{Nome:} \underline{\hspace{3.2cm}} \quad \textbf{Data:} \underline{\hspace{1.6cm}}

\noindent\textbf{\faMusic~O som /dj/ {[dZ]}:} dia, dica, dente, cidade, amigo (variação regional).
\newline\textbf{\faPencilRuler~Quem faz esse som:} D + I (regional). \hfill \textbf{\faCheckCircle~Voz:} com voz (variação regional).

\vspace{0.4em}
\textbf{A) CAÇA AO SOM} --- Em muitas regiões do Brasil a pronúncia varia. Leia devagar e compare: você fala TCHIA ou TIA? DJIA ou DIA?

\begin{center}
\begin{tabular}{@{}ll@{}}
\fbox{\textbf{SIM}} \textit{dia} \quad \fbox{\textbf{NÃO}} \textit{dado} \\
\fbox{\textbf{SIM}} \textit{dica} \quad \fbox{\textbf{NÃO}} \textit{bola} \\
\fbox{\textbf{SIM}} \textit{dente} \quad \fbox{\textbf{NÃO}} \textit{dedo} \\
\fbox{\textbf{SIM}} \textit{cidade} \quad \fbox{\textbf{NÃO}} \textit{casa} \\
\end{tabular}
\end{center}

\textbf{B) SEPARE} --- Escreva as palavras na coluna certa.

\begin{center}
\begin{tabular}{@{}p{4.6cm}|p{4.6cm}@{}}
\hline
\textbf{TEM o som} & \textbf{NÃO tem o som} \\
\hline
\textbf{amigo} & \textbf{sapo} \\
\textbf{dia} & \textbf{pipa} \\
\textbf{dente} & \textbf{teto} \\
\hline
\end{tabular}
\end{center}

\textbf{C) COMPLETE} --- Complete com a grafia do som e confira entre parênteses.

\begin{center}
\large \underline{\hspace{0.7cm}}ia (dia) \quad \underline{\hspace{0.7cm}}ica (dica) \quad \underline{\hspace{0.7cm}}ente (dente) \quad cidad\underline{\hspace{0.7cm}} (cidade)
\end{center}

\textbf{D) OUÇA E COMPARE} --- D seco e D de DIA.

\begin{center}
\begin{tabular}{@{}p{4.2cm}p{6.2cm}p{1.6cm}@{}}
\hline
\textbf{Palavra} & \textbf{Compare} & \textbf{Marque} \\
\hline
\textbf{dado} & D seco [d] & \textbf{NÃO} \\
\textbf{dia} & D + I vira DJ [dj] & \textbf{SIM} \\
\textbf{dedo} & seco & \textbf{NÃO} \\
\hline
\end{tabular}
\end{center}

\caixapausa{Se precisar, pare aqui. Respire fundo e continue quando quiser.}
\caixapista{Língua encosta e solta COM voz (DJ). Mão na garganta: vibra.}
\end{exercicio}

\plancheta{Atividade de Som 38/38 --- H mudo --- O SEGREDO DO H}
\begin{exercicio}[Som H mudo [sem som] --- Atividade de Som (não de letra nem de sílaba)]
\noindent\textbf{Nome:} \underline{\hspace{3.2cm}} \quad \textbf{Data:} \underline{\hspace{1.6cm}}

\noindent\textbf{\faMusic~O som H mudo {sem som}:} hoje (lê-se OJE), hora (ORA), hotel (OTEL), herói (ERÓI).
\newline\textbf{\faPencilRuler~Quem faz esse som:} H (sem som). \hfill \textbf{\faCheckCircle~Voz:} sem som (mudo).

\vspace{0.4em}
\textbf{A) CAÇA AO SOM} --- Leia em voz alta e compare os sons. Marque a coluna certa para cada palavra.

\begin{center}
\begin{tabular}{@{}ll@{}}
\fbox{\textbf{SIM}} \textit{hoje} \quad \fbox{\textbf{NÃO}} \textit{mala} \\
\fbox{\textbf{SIM}} \textit{hora} \quad \fbox{\textbf{NÃO}} \textit{bola} \\
\fbox{\textbf{SIM}} \textit{hotel} \quad \fbox{\textbf{NÃO}} \textit{dedo} \\
\fbox{\textbf{SIM}} \textit{herói} \quad \fbox{\textbf{NÃO}} \textit{casa} \\
\end{tabular}
\end{center}

\textbf{B) SEPARE} --- Escreva as palavras na coluna certa.

\begin{center}
\begin{tabular}{@{}p{4.6cm}|p{4.6cm}@{}}
\hline
\textbf{TEM o som} & \textbf{NÃO tem o som} \\
\hline
\textbf{horta} & \textbf{sapo} \\
\textbf{humor} & \textbf{pipa} \\
\textbf{hino} & \textbf{teto} \\
\hline
\end{tabular}
\end{center}

\textbf{C) COMPLETE} --- Complete com a grafia do som e confira entre parênteses.

\begin{center}
\large \underline{\hspace{0.7cm}}oje (hoje (leia OJE)) \quad \underline{\hspace{0.7cm}}ora (hora (leia ORA)) \quad \underline{\hspace{0.7cm}}otel (hotel (leia OTEL)) \quad \underline{\hspace{0.7cm}}ino (hino (leia INO))
\end{center}

\textbf{D) OUÇA E COMPARE} --- H não tem som.

\begin{center}
\begin{tabular}{@{}p{4.2cm}p{6.2cm}p{1.6cm}@{}}
\hline
\textbf{Palavra} & \textbf{Compare} & \textbf{Marque} \\
\hline
\textbf{hora} & leia sem o H: ORA & \textbf{SIM} \\
\textbf{ora} & palavra de verdade! compare & \textbf{NÃO} \\
\textbf{hoje} & leia sem o H: OJE & \textbf{SIM} \\
\hline
\end{tabular}
\end{center}

\caixapausa{Se precisar, pare aqui. Respire fundo e continue quando quiser.}
\caixapista{O H não tem som: boca neutra, soa a vogal seguinte. Mostre que HOJE se lê OJE (sem sopro).}
\end{exercicio}
